%%%%%%%%%%%%%%%%%%%%%%%%%%%%%%%%%%%%%%%%%%%%%%%%%%%%%%%%%%%%%%%%%%%%%
%%                                                                 %%
%% Please do not use \input{...} to include other tex files.       %%
%% Submit your LaTeX manuscript as one .tex document.              %%
%%                                                                 %%
%% All additional figures and files should be attached             %%
%% separately and not embedded in the \TeX\ document itself.       %%
%%                                                                 %%
%%%%%%%%%%%%%%%%%%%%%%%%%%%%%%%%%%%%%%%%%%%%%%%%%%%%%%%%%%%%%%%%%%%%%


%%======================================================%%
%% to compile with pdflatex/xelatex use pdflatex option %%
%%======================================================%%


\documentclass[12pt]{article}% Math and Physical Sciences Reference Style %% Uncomment this for submission
%\documentclass[pdflatex,sn-mathphys]{sn-jnl-no-heading}% Math and Physical Sciences Reference Style 

%%\documentclass[referee,sn-basic]{sn-jnl}% referee option is meant for double line spacing

%%=======================================================%%
%% to print line numbers in the margin use lineno option %%
%%=======================================================%%

%%\documentclass[lineno,sn-basic]{sn-jnl}% Basic Springer Nature Reference Style/Chemistry Reference Style



%%%% Standard Packages
%%<additional latex packages if required can be included here>
\usepackage[a4paper]{geometry}
\geometry{
 paperheight = 297mm, % Tamaño A4: 297mm
 paperwidth = 210mm, % Tamaño A4: 210mm
 textheight = 247mm, % Valor máximo: 247mm
 textwidth =160mm, % Valor máximo: 160mm
 centering, % Lo habitual
 top = 25mm, % Lo habitual: 25mm
% includehead, % Incluir la cabecera en el área del texto
}
\setlength{\parindent}{7pt}
\usepackage{amssymb}               % para cargar las fuentes de la A.M.S.

\usepackage{amsmath}               % para cargar las fuentes de la A.M.S.
\usepackage{hyperref}   % Autoref

\usepackage{ dsfont }
\usepackage{enumitem}            %autoref enumerate
\usepackage{faktor}




% Shortcuts
\newcommand\N{\mathbb{N}}
\newcommand\Real{\mathbb{R}}
\newcommand\Comp{\mathbb{C}}
\newcommand{\id}[1]{\mathds{1}_{#1}}
\newcommand{\mb}[1]{{\mathbf{#1}}}
\newcommand{\sch}[1]{{\mathcal{S}(\Real^#1)}}
\newcommand{\dsch}[1]{{\mathcal{S}'(\Real^#1)}}
\newcommand{\wickdel}[2]{:#1:\rvert_{\scaleto{#2}{3pt}}}
\newcommand{\braket}[1]{{\langle #1\rangle}}
\newcommand{\bra}[1]{\langle #1\lvert}
\newcommand{\ket}[1]{\rvert #1 \rangle}

%%%%

%\usepackage[normalem]{ulem}

%\jyear{2022}%
%=====================================================%
%% THEOREMS
%%====================================================%
\usepackage{aliascnt}   % second counter for autoref

%\theoremstyle{thmstyleone}%
%%\newtheorem{theorem}{Theorem}%  meant for continuous numbers
\newtheorem{theorem}{Theorem}[section]% meant for section wise numbers
%% optional argument [theorem] produces theorem numbering sequence instead of independent numbers for Proposition
\newaliascnt{proposition}{theorem}
\newtheorem{proposition}[proposition]{Proposition}% 
\newaliascnt{corollary}{theorem}
\newtheorem{corollary}[corollary]{Corollary}% 


%%\newtheorem{proposition}{Proposition}% to get separate numbers for theorem and proposition etc.

%\theoremstyle{thmstyletwo}%

\newtheorem{remark}{Remark}%

%\theoremstyle{thmstylethree}%
\newtheorem{definition}{Definition}%
\newtheorem{example}{Example}% 

\aliascntresetthe{proposition}
\aliascntresetthe{corollary}


\providecommand*\theoremautorefname{{\bf Theorem}}
\providecommand*\propositionautorefname{{\bf Proposition}}
\providecommand*\corollaryautorefname{{\bf Corollary}}
\providecommand*\exampleautorefname{{\it Example}}
\providecommand*\definitionautorefname{{\bf Definition}}

\raggedbottom %The \raggedbottom declaration makes all pages the height of the text on that page. No extra vertical space is added.
%%\unnumbered% uncomment this for unnumbered level heads


%Atuthors
\usepackage{authblk}

\begin{document} 

\title{Geometric flavours of Quantum Field theory on a Cauchy hypersurface. Part I: Geometric quantization }

%%=============================================================%%
%% Prefix	-> \pfx{Dr}
%% GivenName	-> \fnm{Joergen W.}
%% Particle	-> \spfx{van der} -> surname prefix
%% FamilyName	-> \sur{Ploeg}
%% Suffix	-> \sfx{IV}
%% NatureName	-> \tanm{Poet Laureate} -> Title after name
%% Degrees	-> \dgr{MSc, PhD}
%% \author*[1,2]{\pfx{Dr} \fnm{Joergen W.} \spfx{van der} \sur{Ploeg} \sfx{IV} \tanm{Poet Laureate} 
%%                 \dgr{MSc, PhD}}\email{iauthor@gmail.com}
%%=============================================================%%


\author[1,2,3]{{José Luis} {Alonso}}
\author[1,2,3]{{Carlos} {Bouthelier-Madre}}
\author[1,2,3]{{Jesús} {Clemente-Gallardo}}
\author[1,3]{{David} {Martínez-Crespo}}

%\equalcont{These authors contributed equally to this work.}

\affil[1]{{Department of Theoretical Physics},{University of Zaragoza}, {{Facultad de Ciencias, Campus San Francisco},{Zaragoza},{50009},  {Spain}}}

\affil[2]{{Institute for Biocomputation and Physics of Complex Systems}, {University of Zaragoza}, {{Edificio I+D-Campus Río Ebro}, {Zaragoza},{50018},  {Spain}}}

\affil[3]{{Center for Astroparticles and High Energy Physics CAPA}, {University of Zaragoza}, {{Facultad de Ciencias, Campus San Francisco}, {Zaragoza}, {50009},  {Spain}}}

\maketitle



%%==================================%%
%% sample for unstructured abstract %%
%%==================================%%

\abstract{
In this paper we describe the differential geometric properties of a quantum scalar field in terms of Hida-Malliavin calculus. We focus our attention on the role played by the different mathematical structures of the Cauchy hypersurface which constitutes the domain of the classical fields and the geometrical structures that the classical field phase-space is endowed with. We argue that the classical phase space must be modelled over the strong dual of a Nuclear Frechèt space and it is endowed with a complex structure, which allows to define a Kähler structure on it using the canonical symplectic structure. This mathematical structure determines all required ingredients to describe quantum objects through a geometric quantization procedure, analogous to the usual construction on a classical mechanics phase space. We will discuss in detail two different polarizations, one complex and one real; and will compare three different quantizations frequently used in Physics: the usual Schrödinger quantization, the analogous one using the field-momenta instead of the fields; and holomorphic quantization based on creation and annihilation operators. We will see how the necessary ingredients of the quantization procedures take a simple form when written in terms of Malliavin tools.   In the subsequent paper, these tools will be used to characterize the set of quantum states as a Poisson manifold and the unitary dynamics generated by Schrödinger formalism as integral curves of a Hamiltonian vector field. 

\textbf{Keywords: }Quantum fields,  Malliavin calculus, Schrödinger Picture, Holomorphic Picture, Momentum-field Picture, Infinite dimensional Fourier transform, Infinite dimensional Segal-Bargmann Transform, Infinite dimensional geometry}







%%\pacs[JEL Classification]{D8, H51}

%%\pacs[MSC Classification]{35A01, 65L10, 65L12, 65L20, 65L70}



%{\color{red} Hay que quitar el índice del preprint}
\tableofcontents

\section{Introduction}

The mathematical description of Quantum Field Theory on an arbitrary background has always been a complicated task. During the past century, most efforts focused on the algebraic side \cite{Haag1964, Streater1989,Doplicher1969,Doplicher1969a,Doplicher1971,Doplicher1974,Buchholz1982,Buchholz2020}. That approach proved to be very successful in encoding the properties of quantum operators and states in a mathematically consistent framework. Nonetheless, it did not consider other mathematical aspects of the description, as  the differential geometric ones; which were not described in detail either in the other main mathematical approach to QFT, non-commutative geometry (\cite{connesNoncommutativeGeometry1994,connesRenormalizationQuantumField2000}), more focused, again, on algebraic tools.   In particular, the geometrical description of the space of fields justifying, for instance, the existence of a Poisson bracket or of a symplectic form, relays on the definition of the suitable differentiable structure on the fields phase-space, which in this case is cumbersome to build. 

 Furthermore, one of our interests in this construction comes from its potential use in a hybrid quantum-classical model which couples the quantum field and a classical gravitational field. The description of gravity largely depends on differential geometric tools, thus we taylor our  description of  the quantum scalar field  with those tools too. Our main reference for that will be the work  \cite{Kibble1979}, where Kibble defines a Poisson structure on the space of pure states of the Schrödinger-Hamiltonian description of  QFT.   
 In our work we consider the same idea and trace back the origin of different geometric tools  introduced on the phase spaces of fields and pure states. We aim to carefully study the interdependence of each approach to quantization with the complex structure of the classical phase space or on the metric tensor defined on the space-time.  This aim justifies why we prefer  a non-covariant description of the quantum fields (although the global model is covariant) and the choice of the Hamiltonian 3+1 formalism for  gravity.      


Our construction is largely built on the results of Hida-Malliavin calculus developed during the last two decades of the XXth century to study stochastic variational systems and applied with huge success to model financial systems \cite{Hida1980,Hida1993,Nunno2009, Wester2003}. Those tools allow to present the Hamiltonian descriptions of pure quantum states in a simplified language. The set of states has been usually studied with different algebraic or analytical tools under the names of  Schrödinger and holomorphic pictures \cite{Corichi2004,Oeckl2012,Oeckl2012a}.  In this paper, we re-obtain those results but from a geometrical perspective which will allow us to build a description in terms of Poisson structures in the second part of this series.
  This geometrical flavor requires of a careful description of the different structures and of the relations between them required to build the final quantum description.  Hence, we want to build a consistent geometric description of  different set of tools developed in the last years to model QFT.

The structure of the paper is as follows. In Section \ref{sec:math_preliminaries} we stablish notation and summarize the main geometrical tools used in the paper. As we want to do differential geometry, we will briefly summarize how to use the structure of nuclear spaces  to build a differentiable structure on the set of fields. Furthermore, we will discuss  the use of rigged Hilbert spaces to identify these nuclear spaces with suitable subspaces of their duals to be able to map vectors and covectors on those field manifolds. This mechanism provides a practical framework to deal with tensor calculus.  A very important ingredient  that justifies the use of nuclear spaces to model the space will be  integral calculus. In particular the origin, existence and properties of the set of gaussian measures on the space of fields, which are required to build the Hilbert space structure on the set of quantum pure states. Finally, we will introduce tools to relate our framework with the usual approach to QFT. In particular, we will see how the Segal isomorphism allows us to relate the quantum states with the Wiener-Ito chaos decomposition which defines an isomorphism with the (bosonic) Fock space which is most commonly used in QFT. We will also discuss the Seagal-Bargmann transform which will be used later in the paper to relate the real and holomorphic description of the quantum fields. And finally, we will introduce the main tools of Malliavin calculus, in particular the Malliavin derivative and the Skorokhod integral; and discuss in some detail the relation of these objects with the creation and annihilation operators in QFT.  The rest of the paper describes step-by-step how can we build the Hilbert space of the space of quantum fields, its properties of and how they intertwine with the differential geometric properties of the phase space of classical fields.  Thus, 
Section \ref{sec:CFTQuantizar} analyzes the geometrical structure of the classical phase space of fields and how the introduction of a complex structure in that set determines the structure of the Hilbert space of quantum fields. This issue will be crucial when constructing the hybrid quantum-classical model.  Sections \ref{sec:QFTCauchy} and \ref{subsec:holom} present the geometric quantization of the resulting classical field theory, focusing on the role of the different polarizations and how the different tools introduced allow us to write the isomorphisms relating them at the level of linear operators. In the  fourth section we analyze geometric quantization in general, while the fifth considers in some detail three important special cases and their relations: Holomorphic, Schrödinger  and momentum-field quantization. To relate the three of them we must introduce integral transforms generalizing Fourier and Segal-Bargmann transforms to an infinite dimensional domain with a nontrivial complex structure. Section \ref{sec:orderings} studies in some detail the extension of that construction to higher order operators, discussing Weyl quantization and its relation with Wick quantization. Finally, Section \ref{sec:conclusions} summarizes the main results of the paper and describe their applications in its second part.


\section{Mathematical preliminaries}
\label{sec:math_preliminaries}
In the following sections we will be constructing a differential-geometrical description of the phase spaces of classical and quantum field theories, which are, in both cases, infinite dimensional linear spaces. In order to do that, we will require modelling those manifolds on topological vector spaces where we can construct a differential calculus. Besides, in the case of the quantum phase space, we will also require the operators representing our physical magnitudes to be essentially self-adjoint on their domains, which will be submanifolds of the total space.  In some occasions, we may use a Hilbert space as modelling space, but this is not the general case. The most general framework is provided by dense subsets of functions, with enough properties to admit a well-behaved notion of continuity and differentiability. The best suited sets of such functions will be Nuclear-Frechèt (NF) spaces, whose main properties are presented below, with a few relevant examples. Furthermore, we will require to be able to define rigorously Borel probability measures on those manifolds, for the definition of quantum states to make sense so we are forced to consider also strong duals of NF spaces (DNF) as models of geometry. As gaussian measures are the most commonly used in QFT, this will be our preferred choice. 

Furthermore, the construction of tensors on these manifold will require considering both tangent and cotangent spaces. As the dual spaces of the nuclear spaces are spaces of distributions, tangent and cotangent spaces will be  non-isomorphic. In order to define a practical model, we will consider the notion of Gel'fand triple to be able to model vectors and covectors as elements of a suitable Hilbert space, which will provide us with suitable tensors on the tangent or on suitable subsets of the cotangent space. 

Finally, we will analyze in detail the nuclear space used to model the set of quantum fields, the Hida test functions, and the corresponding Gel'fand triple. It is just one example of nuclear space, but its importance for the rest of the paper requires a specific analysis. 

%\subsection{Nuclear spaces and Gel'fand triples}

\subsection{Nuclear spaces: doing differential geometry in infinite dimensions}
It is a well known fact that, in order to define gaussian Borel probability measures in infinite dimensions, the natural setting is that of Frechèt Nuclear spaces (FN) and its dual spaces (DFN). This is discussed thoroughly by \cite{Gellfand}. In this section we will summarize the main ingredients of integration theory on the dual pair  $(\mathcal{N},\mathcal{N}')$ where $\mathcal{N}$ is a FN space and $\mathcal{N}'$ is its dual endowed with the strong topology\footnote{In this case we endow the DFN space with the strong topology, which coincides with the Mackey topology, but the same Borel $\sigma$-algebra is generated from the weak topology. }. In particular, we are interested in
$\big(C^{\infty}_c(K),D'(K)\big)$, the compactly supported functions and distributions over a compact subset $K\subset \mathbb{R}^{d}$ and $\big(\mathcal{S}(\mathbb{R}^d),\mathcal{S}'(\mathbb{R}^d)\big)$ the space of functions of fast decay and tempered distributions. 
 
For completeness we define here the functional spaces to be used in a way that fits our purpose.

\begin{definition}
    \label{def:Nuclear}
    {\bf (Nuclear Space)} A nuclear space $X$ is a  topological vector space whose topology is given by a family of Hilbert seminorms  such that, if $X_{p}$ is the Hilbert space obtained as  the completion of the preHilbert  space $(X/p^{-1}(0), p)$, then there is a larger seminorm $q$ such that $X_q\hookrightarrow X_p$ is a Hilbert-Schmidt operator. 
\end{definition}

For infinite dimensional topological vector spaces $X,Y$ tensor products  can be endowed with a family of  different topologies between the projective $X\otimes_\pi Y$ and injective $X\otimes_\epsilon Y$ tensor product induced topologies.  If one or both spaces are nuclear both topologies coincide and therefore there is only one notion of tensor product $X\otimes Y$. We will  use this fact in following sections to define polynomials over $X$.

\begin{definition}
    \label{def:FN}
    {\bf (FN Frechèt Nuclear Space)} A Nuclear space which is also Frechèt  $\mathcal{N}$  can be regarded as the projective limit of a countable family of separable Hilbert spaces $\{\mathcal{H}_p\}_{p=0}^\infty$, with scalar products $\langle \cdot, \cdot  \rangle_p$ \ such that for $q>p$     
    the injection $\mathcal{H}_q\hookrightarrow \mathcal{H}_p$ is Hilbert-Schmidt. Thus
    $$
    \mathcal{N}=\bigcap_{p=0}^\infty\mathcal{H}_p 
    $$
    with the projective limit topology. 
\end{definition}

 Riesz representation theorem allows us to identify a Hilbert space with its dual. By convention we set $\mathcal{H}_0'=\mathcal{H}_0^*$ where $*$ represents complex conjugation. This implies that in the inclusion $\mathcal{H}_p\hookrightarrow\mathcal{H}_0$, for $p>0$, the space $\mathcal{H}_p$ can no longer be identified with its dual. We will denote this dual space by $\mathcal{H}_0\hookrightarrow\mathcal{H}_{-p}$. The next definition can be regarded as a consequence of the last one. 

\begin{definition}
    \label{def:DFN}
    {\bf (DFN Dual of a Frechèt Nuclear Space)} The strong dual of a FN space  $\mathcal{N}'$  can be regarded as the inductive limit of a countable family of Hilbert spaces $\{\mathcal{H}_{-p}\}_{p=0}^\infty$ such that for $q>p$
    the injection $\mathcal{H}_{-q}\hookleftarrow \mathcal{H}_{-p}$ is the transpose of a  Hilbert-Schmidt operator. Thus
    $$
    \mathcal{N}'=\bigcup_{p=0}^\infty\mathcal{H}_{-p} 
    $$
    with the injective limit topology. A DFN space is always nuclear. Furthermore it is a Nuclear Silva space \cite{Kriegl1997}.
\end{definition}




\subsection{Rigged Hilbert spaces and main examples}

With the definitions given so far we notice that the space $\mathcal{H}_0$ plays an important role. This motivates the consideration of a triple of spaces therefore leading to the concept of rigged Hilbert space. 

\begin{definition} {\bf (Rigged Hilbert Space)}
    \label{def:RHS}
    Assume that $\mathcal{H}_0$ is a separable Hilbert space then, then we define a rigged Hilbert space or Gel'fand triple as the triple of spaces
\begin{equation*}
\label{eq:rigged}
\mathcal{N}\subset \mathcal{H}_0 \subset \mathcal{N}'
\end{equation*}
where $\mathcal{N}$ is a FN space and $\mathcal{N}'$ its dual and each inclusion is continuous. 

\end{definition}

In the structure here presented $\mathcal{H}_0$ is said to realize the duality, which means that the scalar product of $\mathcal{H}_0$ identifies an element of $\mathcal{N}$ with its dual in $\mathcal{N}'$. 
$\mathcal{N}$ is a dense subspace of the Hilbert space.
This must be seen as a way of parametrizing  $\mathcal{N}$ as a subset of $\mathcal{N}'$, in general different Hilbert spaces can be used to parametrize the duality of the same pair $(\mathcal{N},\mathcal{N}')$.



In the rigging there are three main ingredients. Firstly,  elements  $\xi_{\mathbf{x}}\in\mathcal{N}$ will be typically well defined functions (see examples below) and will be marked by a bold subindex.  Secondly,  elements $\varphi^{\mathbf{x}}\in\mathcal{N}'$ will be distributions with no pointwise meaning in general and they will be marked by a bold superindex. Sometimes we will denote the value of the function at a point $x$  with  a regular subindex $\xi_x$, this can't be done in general for distributions. With those two ingredients the duality is denoted, mimicking Einstein summation convention, 
\begin{equation}
\label{eq:dualityNotationDistributions}
\xi_{\mathbf{x}}\in\mathcal{N}, \varphi^{\mathbf{x}}\in\mathcal{N}' \textrm{ paired by } \xi_x\varphi^x :=\langle\varphi^{\mathbf{x}},\xi_{\mathbf{x}}\rangle \in \mathbb{R}\ \ (\mathbb{C}).
\end{equation}

The third and last ingredient is the dual pairing of $\mathcal{H}_0$. This element can be seen in a threefold way
\begin{enumerate}[label=\alph*)]
    \item Dual pairing between elements in $\mathcal{N}\subset \mathcal{H}_0$ with $\langle\xi^1_{\mathbf{x}} ,\xi^2_{\mathbf{y}}\rangle =\xi^1_x\Delta^{xy}\xi_y^2$
    \item Dual pairing between elements in  $ \mathcal{H}_0\subset \mathcal{N}'$ with $\langle\varphi_1^{\mathbf{x}} ,\varphi_2^{\mathbf{y}}\rangle =\varphi_1^xK_{xy}\varphi^y_2$ 
    \item Injection of ${\Delta^{\mathbf{xy}}}:\mathcal{N}\hookrightarrow \mathcal{N}'$ trough the pairing of $\mathcal{H}_0$ such that $\xi_x\Delta^{xy}=\xi^y$ with a densely defined inverse $\varphi^xK_{xy}=\varphi_y$ such that $\Delta^{xz}K_{zy}=\delta^x_y$ is the evaluation mapping.
\end{enumerate}


The main examples of FN spaces, and the spaces that play the roles of model spaces for the classical field theories in this work are the next two.

\begin{example}\label{ex:testFunctions}{\bf (Smooth functions of compact support on a compact $K$)}
    Let $K\subset \mathbb{R}^d$ then the space $C^{\infty}_c(K)$ is FN with the initial topology induced from the family of Hilbertian norms

\begin{equation}
\label{eq:hilbertianNormsCompactSupport}
\lVert\varphi\rVert_{K,n}^2=\sum_{\lvert\alpha\rvert\leq n}\int_{K}\lvert\partial^\alpha\varphi(x)\rvert^2 dx^d,
\end{equation}
where $\alpha\in \mathbb{N}_0^d$ is a multi-index and $\lvert\alpha\rvert=\alpha_1+\cdots+\alpha_d$.  Its dual is the DFN space $D'(K)$. 

Notice that we could also define $C^\infty_c(U)$, for $U\subset \mathbb{R}$ an open set. This is a nuclear  but not  a Frechèt space; instead it is a Nuclear LF space (see \cite{Gellfand} for details). For our purposes we will limit our attention to the compact domain in order to reduce the technicalities.
\end{example}



\begin{example}{\bf (Smooth functions of fast decay)}\label{ex:schwartFunctions} The space $\mathcal{S}(\mathbb{R}^d)$ of smooth functions of fast decay is FN with the initial topology induced from the family of Hilbertian norms

    \begin{equation}
        \label{eq:hilbertianNormsSchwartz}
        \lVert\varphi\rVert_{m,n}^2=\sum_{\lvert\alpha\rvert\leq n, \lvert\beta\rvert\leq m}\int_{\mathbb{R}^d}\prod_{i=1}^d(1+x_i^2)^{2\beta_i}\lvert \partial^\alpha\varphi(x)\rvert^2 dx^d
    \end{equation}
    where $\alpha,\beta\in \mathbb{N}_0^d$ are multi-indices. This topology is equivalent to the one usually introduced via supremum type seminorms. Its dual $\mathcal{S}'(\mathbb{R}^d)$ is the DFN space of tempered distributions. This space is the one commonly used in white noise analysis \cite{Hida1993}  
\end{example}


\begin{example}\label{ex:testFunctionsManifold}{\bf (Smooth functions of compact support on a compact Manifold )} For $\Sigma$ a compact manifold the definition of $C^\infty_c(\Sigma)$ goes on as in \autoref{ex:testFunctions}. If we endow the manifold with a riemannian metric $h$ there is a canonical volume form 
\begin{equation}
\label{eq:riemannianVolumeform}
\int_\Sigma f\mathrm{dVol}_h=\int_U dx^d \sqrt{\lvert h(x)\rvert} f(x),
\end{equation}
for $f$ supported on a subset of $\Sigma$ covered by  a single chart with coordinates in $U$. This, in turn, leads to a canonical Hilbert space of square integrable functions $L^2(\Sigma,\mathrm{dVol}_h)$ and therefore to the rigged Hilbert space 
\begin{equation}
\label{eq:riggedOnAManifold}
C^\infty_c(\Sigma)\subset L^2(\Sigma,\mathrm{dVol}_h) \subset D'(\Sigma)
\end{equation}  
In this picture $D'(\Sigma)$ are densities that are used to integrate the elements of $C^\infty_c(\Sigma)$ and this nontrivial structure is reframed in the rigged Hilbert space with a nontrivial representation of the scalar product as an integral kernel.  

As discussed in \autoref{def:RHS} the scalar product on a rigged Hilbert space can be considered in a threefold way, as we can illustrate over a chart with coordinates in $U\subset \mathbb{R}^d$

\begin{enumerate}[label=\alph*)]
        \item Dual pairing between elements in $C^\infty_c(\Sigma)\subset L^2(\Sigma,\mathrm{dVol}_h)$ with $$\langle f_\mb{x} ,g_\mb{y}\rangle =f_x\delta^{xy}g_y=\int_U dx^d\sqrt{\lvert h(x)\rvert} f(x)g(x)$$
    \item Dual pairing between elements in  $L^2(\Sigma,\mathrm{dVol}_h) \subset D'(\Sigma)$ with $$\langle\phi^{\mathbf{x}} ,\varphi^{\mathbf{y}}\rangle =\phi^x\delta_{xy}\varphi^y=\int_U\frac{dx^d}{\sqrt{\lvert h(x)\rvert}} \phi(x)\varphi(x),$$
    where $\phi(x),\varphi(x)$ are densities. Notice that the factor $\sqrt{\lvert h(x)\rvert}$ appears in the denominator and $\frac{\phi(x)}{\sqrt{\lvert h(x)\rvert}}, \frac{\varphi(x)}{\sqrt{\lvert h(x)\rvert}}$ are $L^2(\Sigma, \mathrm{dVol}_h)$ representatives.
    \item We can also see the injection ${\delta^{\mathbf{xy}}}:C^\infty_c(\Sigma) \hookrightarrow D'(\Sigma)$  such that
    $$f^y=f_x\delta^{xy}=\int_U dx^d\sqrt{\lvert h(x)\rvert}\delta^d(x-y) f(x)=\sqrt{\lvert h(y)\rvert}f(y)$$
    is a representative of a density. On the other hand $\delta_{\mathbf{xy}}$ is a densely defined inverse 
    $$\varphi^y=\varphi^x\delta_{xy}=\int_U \frac{dx^d}{\sqrt{\lvert h(x)\rvert}}\delta^d(x-y) \varphi(x)=\frac{\varphi(y)}{\sqrt{\lvert h(y)\rvert}}$$
    such that $\delta^{xz}\delta_{zy}=\delta^x_y$ is the evaluation mapping
    $$\delta^x_yf_x=f_y=f(y)$$
    which only has a definite sense over $C^\infty_c(\Sigma)$.
\end{enumerate}

The differences between $\delta^{\mathbf{xy}}$, $\delta_{\mathbf{xy}}$ and $\delta^{\mathbf{x}}_{\mathbf{y}}$ are often misinterpreted since the three of them can be identified with the Dirac delta in $\mathbb{R}^d$ with the Lebesgue measure. 
\end{example}

\subsection{Defining a scalar product: Gaussian measures}
As we are intending to create a quantum description of our physical state, we must endow the space of quantum states with a scalar product.  As we mentioned above, the most commonly used measures in QFT correspond to gaussian measures, often referred to a given vacuum state. 
We will review below their main definition and  properties, first for the simple case of real test functions, and then in the case of complex ones, which will be more useful for us later.
\subsubsection{The real case}
It is easy to define a Gaussian measure over $\mathcal{N}'$. Consider the $\sigma$-algebra generated by its cylinder subsets. It corresponds to the Borel $\sigma$-algebra of the weak and strong topologies of $\mathcal{N}'$.  Once the characteristic functional,
\begin{equation}
    \label{eq:chartacteristicFunctional}
        C(\xi_{\mathbf{x}})=\int_{\mathcal{N}'}D\mu(\varphi^{\mathbf{x}}) e^{\xi_x\varphi^x} =e^{-\frac{1}{2}{\xi}_x\Delta^{xy}\xi_y},
\end{equation}
is known, Minlos theorem \cite{Gellfand,Hida1993} ensures that, if ${\xi}_x\Delta^{xy}\xi_y$ is continuous with respect to the topology of $\mathcal{N}$, there is a unique Borel probability measure that fulfills \eqref{eq:chartacteristicFunctional}. Considered with respect to a Gel'fand triple, the bilinear symmetric form $\Delta$ corresponds to the representation of the scalar product of the Hilbert space on the set of test functions and defines the covariance of the gaussian.


\subsubsection{The complex case}
If we consider the case of complex test functions $\mathcal{N}_\mathbb{C}$, of any of the cases above, most of the properties work in the same form.  Let us consider the complexification of the space of distributions such that  an element $\phi^{\mathbf{x}}\in\mathcal{N}'_\mathbb{C}$ is written as $\phi^{\mathbf{x}}=\frac{1}{\sqrt{2}}(\varphi^{\mathbf{x}}+i\pi^{\mathbf{x}})$ for $(\varphi^{\mathbf{x}},\pi^\mathbf{x})\in \mathcal{N}'\times \mathcal{N}'$. Complex test functions can also be written in complex coordinates, for $\rho_{\mathbf{x}}\in\mathcal{N}_\mathbb{C}$ it follows  $\rho_x=\frac{1}{\sqrt{2}}(\xi^{\varphi}_x+i\xi^{\pi}_x)$ and $\bar\rho_x=\frac{1}{\sqrt{2}}(\xi^{\varphi}_x-i\xi^{\pi}_x)$. If we endow each copy of the cartesian product with the measure given by \eqref{eq:chartacteristicFunctional},  we can define a measure to the complex domain  defined by the characteristic functional  

\begin{equation}
\label{eq:complexCharacteristicFunctional}
C_c(\rho_{\mathbf{x}},\bar\rho_{\mathbf{y}})=\int_{\mathcal{N}_{\mathbb{C}}'}D\mu_c(\bar \phi^{\mathbf{x}}, \phi^{\mathbf{x}}) e^{i({\rho_x\bar{\phi}^x+\bar\rho_x{\phi}^x})} =e^{-\bar{\rho}_x\Delta^{xy}\rho_y},
\end{equation}
where now the covariance $\Delta^{xy}$ will be a hermitian form on $\mathcal{N}_\mathbb{C}$, i.e. $\Delta^{xy}=(\Delta^*)^{yx}$. 
Notice that the complex measure is the product measure of two real ones

\subsubsection{Integral tranforms on gaussian measures}
Among the different integral transforms defined over a gaussian measure the most relevant for Quantum Field Theory is the $S$ transform. Throughout all the section we will consider a gaussian measure defined on the dual of a nuclear space $\mathcal{N}$ with a Gel'fand triple
$$
\mathcal{N}\subset \mathcal{H}_\Delta \subset \mathcal{N}',
$$
where $\mathcal{H}_\Delta$ represents the Cameron-Martin Hilbert space that is obtained by completion with the covariance of a gaussian measure $D\mu$ with associated characteristic functional \eqref{eq:complexCharacteristicFunctional}. For the real case consider functions defined on the space of distributions $\mathcal{N}'$ then :

\begin{definition}
    \label{def:Stransform}{\bf ($S_\mu$ transform)} For the complex gaussian measure of \eqref{eq:chartacteristicFunctional} we define the $S_\mu$ transform as a generalization of the Laplace transform. For a given $\Psi:\mathcal{N}'\to \mathbb{C}$ it takes the form 
    \begin{equation}
        \label{eq:stransform}
        S_\mu[\Psi](\xi_\mathbf{x})=\int_{\mathcal{N}'}D\mu(\varphi^{\mathbf{x}})\Psi(\varphi^{\mathbf{x}})\exp{\big(\xi_x\varphi^x-\frac12\xi_x\Delta^{xy}\xi_y\big)}
\end{equation}
\end{definition}
This is to be thought as the action of a translation $\varphi^{\mathbf{x}}\to\varphi^{\mathbf{x}}+\Delta^{\mathbf{x}y}\xi_y$ on the domain of integration . For the complex case it is readily defined in (anti)-holomorphic coordinates as 

\begin{definition}
    \label{def:SComplextransform}{\bf ($S_{\mu_c}$ complex transform)} For the complex gaussian measure of \eqref{eq:complexCharacteristicFunctional} and a given $\Psi:\mathcal{N}_{\mathbb{C}}'\to \mathbb{C}$, the transform takes the form 
    \begin{equation}
        \label{eq:stransformcomplex}
        S_{\mu_c}[\Psi](\rho_\mathbf{x},\bar\rho_\mathbf{x})=\int_{\mathcal{N}'_{\mathbb{C}}}D\mu_c(\phi^{\mathbf{x}})\Psi(\phi^{\mathbf{x}},\bar\phi^{\mathbf{x}})\exp{\big({\rho}_x\bar\phi^x+\bar{\rho}_x{\phi}^x-\bar\rho_x\Delta^{xy}\rho_y\big)}
        \end{equation}
\end{definition}

Notice that this definition is nothing but the combination of two copies of \autoref{def:Stransform}, one for each variable, and the required change to holomorphic coordinates.

In the next section we will define relevant physical concepts based on these tools.

\subsection{Modelling quantum field pure states: Wiener Ito decomposition}
In this section we will consider the mathematical tools required to represent the concept of pure state of a quantum field and its main properties, in particular the relation with the concept of physical particles. In this context, the notion of particle in QFT is defined as the eigenstates of the  particle number operator on a Fock space. On the other hand, a QFT pure state, for the scalar field, must be considered as a function of the Hilbert space $L^2(\mathcal{N}'_\mathbb{C}, D\mu_c)$, i.e. is to be thought as a function with domain in the DFN space  $\mathcal{N}'_\mathbb{C}$. The precise relation between the Fock space and $L^2(\mathcal{N}'_\mathbb{C}, D\mu_c)$ is given by the Wiener-Ito decomposition in homogeneous chaos, which we discuss briefly below. In the same lines, we can also identify common tools of gaussian analysis as Malliavin derivatives and Skorokhod integrals with creation and annihilation operators defining in the process the notion of charge and particle number. 



\subsubsection{The Segal isomorphism: Wiener-Ito chaos decomposition and the Fock space}
The classical phase space $\mathcal{M}_C$ can be considered as the set of complex test functions $\mathcal{N}_\mathbb{C}$ once we introduce a suitable complex structure on the set $\mathcal{N}\times \mathcal{N}$ (containing the fields and their momenta). We will start dealing with the complex case and then particularize our result to the real case. If the Gel'fand triple discussed above is considered with respect to complex spaces, we can build the triple

\begin{equation}
\label{eq:complexgelfand}
\mathcal{M}_C\sim \mathcal{N}_\mathbb{C}\subset \mathcal{H}_\Delta\subset \mathcal{N}'_\mathbb{C},
\end{equation}
 where $\mathcal{H}_\Delta$ represents the, now complex, Cameron-Martin Hilbert space. By using this triple, we can consider $\mathcal{M}_C$ as a linear subspace of $\mathcal{N}'_\mathbb{C}$. This means that we will think at our classical fields and their momenta as distributions on the nuclear space of complex smooth functions of compact support on $\Sigma$.

Now, we move to the level of second quantization.  In the next subsection we will consider  test functions defined on the set of classical fields extended by its Gel'fand triple (i.e., on $\mathcal{M}_C \subset \mathcal{N}'_\mathbb{C}$).   Those functions will represent a dense subspace of the Hilbert space of quantum pure states $L^2(\mathcal{N}'_\mathbb{C}, D\mu_c)$ and will allow us to build a differential geometric structure for it.  For now, we will define a basis of orthogonal polynomials and the corresponding system of coordinates for this Hilbert space. We will see that this construction allows us to define an isomorphism with the (bosonic) Fock space constructed from the  (Cameron-Martin) Hilbert space $\mathcal{H}_\Delta$. This isomorphism will allow us to define the set of Hida test functions on $\mathcal{N}'_\mathbb{C}$ in the next section.

As our Hilbert space is $\mathcal{H}=L^2(\mathcal{N}'_\mathbb{C}, D\mu_c)$, we define the pointwise defined polynomials $\mathcal{P}(\mathcal{N}'_\mathbb{C})$, with coefficients in $\mathcal{N}_{\mathbb{C}}^{\otimes n}$ for each monomial of degree $n$. This is, a polynomial of degrees $(n,m)$ is  


\begin{equation}
    \label{eq:polynomialscompl} 
    P_{nm}(\phi^{\mathbf{x}},\bar\phi^{\mathbf{x}})= \sum_{k=0}^n\sum_{l=0}^m C^{kl}_{\vec{x}_k,\vec{y}_l}(\phi^k)^{\vec{x}_k}(\bar\phi^l)^{\vec{y}_l}
\end{equation}
With $C^{kl}_{\vec{\mathbf{x}}_k,\vec{\mathbf{y}}_l}\in \mathcal{N}_{\mathbb{C}}^{\otimes k+l}$ and $(\phi^k)^{\vec{\mathbf{x}}_k}\in {\mathcal{N}_\mathbb{C}'}^{\otimes k}$, $(\bar\phi^l)^{\vec{\mathbf{x}}_l}\in {\mathcal{{N}_\mathbb{C}'}^*}^{\otimes l}$ are the tensor product of copies of the same distribution. Notice that those tensor products are unambiguously defined due to $\mathcal{N}_{\mathbb{C}}$ being an FN space.  This space $\mathcal{P}(\mathcal{N}'_\mathbb{C})$ can be shown to be dense in $\mathcal{H}$ \cite{Hida1993}.

Then we can consider the spaces
\begin{equation}
\label{eq:wickspaces}
\mathcal{W}_\mu^{:(n,m):}=\overline{\mathcal{P}^{(n,m)}}\cap \overline{\mathcal{P}^{{(n-1,m)}}}^\perp\cap \overline{\mathcal{P}^{{(n,m-1)}}}^\perp,
\end{equation}
where $\overline{\mathcal{P}^{(n,m)}}$ represents the completion (in $L^2(\mathcal{N}'_\mathbb{C}, D\mu_c)$) of the set of polynomials of degrees $(n,m)$  and $\overline{\mathcal{P}^{(n-1,m)}}^\perp$ the orthogonal complement of the completion of the set of polynomials of degrees $(n-1,m)$. The spaces  $\mathcal{W}_\mu^{:(n,m):}$  are called homogeneous chaos of degrees $(n,m)$ and are generated by (anti)-holomorphic Wick monomials of the gaussian measure $\mu$ (or, equivalently, of the covariance $\Delta$). They contain a basis of orthogonal polynomials for the Hilbert space $L^2(\mathcal{N}'_\mathbb{C}, D\mu_c)$, generalizing the definition of complex Hermite polynomials to this more abstract setting \cite{Hida1980}.


\begin{definition}
    \label{def:ComplexWickMonomials}{\bf (Complex Wick monomials)} 
    We denote $\mathcal{W}_\Delta(\phi^n\bar{\phi}^m)^{\vec{\mathbf{x}}\vec{\mathbf{y}}}$, with $\vec{x}=(x_i)_{i=1}^n, \vec{y}=(y_i)_{i=1}^m$ and call  complex Wick monomial of degrees $(n,m)$ with respect to the covariance $\Delta^{\mathbf{xy}}$, the distribution belonging to the symmetric tensor product of $n$ copies on $\mathcal{N}'_\mathbb{C}$ and the symmetric product of $m$ copies on ${\mathcal{N}'_\mathbb{C}}^*$ i.e. $\mathcal{W}_\Delta(\phi^n\bar{\phi}^m)^{\vec{\mathbf{x}}\vec{\mathbf{y}}}\in{\mathcal{N}'_\mathbb{C}}^{\hat{\otimes}n}\otimes {\mathcal{N}'_\mathbb{C}}^{*\hat{\otimes}m}$ such that
    \begin{equation}
    \label{eq:ComplexWienerIto}
    \mathcal{W}_\Delta(\phi^n\bar{\phi}^m)^{\vec{x}\vec{y}}=\prod_{i=1}^n\prod_{j=1}^m\frac{\delta}{\delta \bar\rho_{x_i}}\frac{\delta}{\delta \rho_{y_j}}\exp\left[{\bar\rho_z\phi^z+\rho_z\bar\phi^z-\bar\rho_u\Delta^{uv}\rho_v}\right]\bigg\lvert_{\rho_{\mathbf{x}}=0}
    \end{equation}
\end{definition}



Where $L^2_{Hol}$ and $L^2_{\overline{Hol}}$ represent holomorphic (resp. antiholomorphic) square integrable functions, and $\bar{\otimes}$ represents the topological tensor product completed with respect to the natural Hilbert space topology of the tensor product.

 With this set of polynomials, we can write any state in the Hilbert space $L^2(\mathcal{N}_\mathbb{C}, D\mu_c) $ as a linear combination 
    \begin{equation}
    \label{eq:ComplexWienerItodec}
    \Psi(\phi^{\mathbf{x}},\bar\phi^\mathbf{y})=\sum_{n,m=0}^\infty\mathcal{W}_\Delta(\phi^n\bar{\phi}^m)^{\vec{x}_n\vec{y}_m}\psi^{(n,\bar{m})}_{\vec{x}_n\vec{y}_m} 
    \end{equation}
    We call this decomposition Wiener-Ito chaos decomposition in homogeneous chaos.
    Notice that the set of polynomials can be decomposed in holomorphic and anti-holomorphic parts 
    \begin{equation}
    \label{eq:ComplexPolinomials}
    \mathcal{P}_{Hol}(\mathcal{N}'_{\mathbb{C}})\otimes \mathcal{P}_{\overline{Hol}}({\mathcal{N}_{\mathbb{C}}'}^*)\subset L^2\big({\mathcal{N}_{\mathbb{C}}'},D\mu_c\big)= L^2_{Hol}\big({\mathcal{N}_{\mathbb{C}}'}, D\mu_c\big)\bar{\otimes} L^2_{\overline{Hol}}\big({\mathcal{N}_{\mathbb{C}}'}, D\mu_c\big)
    \end{equation}
    where $\bar{\otimes}$ represents the topological tensor product completed with the natural Hilbert topology induced by the tensor product. In terms of the coefficients on the Wiener-Ito decomposition we can write the scalar product as:
    \begin{align}
    \label{eq:complexchaosDecomposition}
    \int_{\mathcal{N}'_\mathbb{C}}D\mu_c(\phi^\mathbf{x})\overline{\Phi}(\phi^\mathbf{x})\Psi(\phi^\mathbf{x})
    &= \nonumber \\ \sum_{n,m=0}^\infty n!m!&\ \overline{\varphi^{(n,\bar{m})}}_{\vec{u}_n\vec{v}_m}\big[(\Delta^*)^n\big]^{\vec{u}_n\vec{x}_n} \big[\Delta^m\big]^{\vec{v}_m\vec{y}_m} \psi^{(n,\bar{m})}_{\vec{x}_n\vec{y}_m}
    =\nonumber \\ &
    \sum_{n,m=0}^\infty n!m! \big\langle\varphi^{(n,\bar{m})}_{\vec{\mathbf{x}}_n\vec{\mathbf{y}}_m},\psi^{(n,\bar{m})}_{\vec{\mathbf{x}}_n\vec{\mathbf{y}}_m}\big\rangle_{\mathcal{H}^{(n,\bar{m})}_\Delta}
    \end{align}

    Therefore weighting coefficients   of the Wiener-Ito decomposition with $\sqrt{n!m!}$ they belong to the Bosonic Fock space $\Gamma_{hol}\mathcal{H}_\Delta\otimes \Gamma_{\overline{hol}}\mathcal{H}_\Delta$.  To find the coefficients of the decomposition is useful to use the $S_{\mu}$ transform as
    \begin{equation}
    \label{eq:wienerItoSeagalwithTtransform}
    \psi^{(n,\bar{m})}_{\vec{x}_n\vec{y}_m}=\frac{1}{n!m!}\prod_{i=1}^n\prod_{j=1}^m K_{x_iu_i}K^*_{y_jv_j}\frac{\partial}{\partial \rho_{u_i}}\frac{\partial}{\partial \bar\rho_{v_j}}  {S_\mu[\Psi](\rho_\mathbf{x},\bar{\rho}_\mathbf{x})} \bigg\lvert_{\rho_\mathbf{x}=0}
    \end{equation}
    Recall that $K_{\mathbf{xy}}$ is the  (weak) inverse of the covariance $\Delta^{\mathbf{xy}}$.
     We see that the decomposition of \eqref{eq:complexchaosDecomposition} is such that
     \begin{equation}
        \label{eq:holantiholdec}
        L^2\big({\mathcal{N}_{\mathbb{C}}'}, D\mu_c\big)=L^2_{Hol}\big({\mathcal{N}_{\mathbb{C}}'}, D\mu_c\big)\bar{\otimes} L^2_{\overline{Hol}}\big({\mathcal{N}_{\mathbb{C}}'}, D\mu_c\big)=\sum_{n=0}^\infty\oplus \left( \sum_{m=0}^n\oplus \mathcal{H}_\Delta^{(n-m,\bar{m})} \right).
    \end{equation}

    This leads us to build an isomorphism between the Hilbert space and the Fock space which associates each state in $L^2({\mathcal{N}_{\mathbb{C}}'}, D\mu_c)$ and its coordinates in the linear decomposition:

\begin{definition}
    \label{def:SegalIso}{\bf (Segal isomorphism)} 
    We will call Segal isomorphism to the mapping


    \begin{equation}
\label{eq:Segal}
\begin{array}{cccc}
    \mathcal{I} :&  L^2\big({\mathcal{N}_{\mathbb{C}}'}, D\mu_c\big) & \longrightarrow &  \Gamma_{hol}\mathcal{H}_\Delta\otimes \Gamma_{\overline{hol}}\mathcal{H}_\Delta
 \\ 
 &  \Psi &\longmapsto &  \mathcal{I}(\Psi)=\left (\sqrt{n!m!}\psi^{(n,\bar{m})}_{\vec x_n \vec y_m} \right )_{n,m=0}^\infty 
\end{array}
\end{equation}
which defines a unitary isomorphism of Hilbert spaces considering the scalar product defined in \eqref{eq:complexchaosDecomposition}. In this decomposition the Fock space is further decomposed into the orthogonal direct sum $\Gamma_{hol}\mathcal{H}_\Delta\otimes \Gamma_{\overline{hol}}\mathcal{H}_\Delta
=\sum_{n=0}^\infty\oplus \left( \sum_{m=0}^n\oplus \mathcal{H}_\Delta^{(n-m,\bar{m})} \right)$.
\end{definition} 
    It is interesting to notice that  $\mathcal{W}_\Delta(\phi^n\bar{\phi}^0)^{\vec{\mathbf{x}}}= (\phi^n)^{\vec{\mathbf{x}}}$. This fact simplifies remarkably computations in $L^2_{Hol}\big({\mathcal{N}_{\mathbb{C}}'}, D\mu_c\big)$ because holomorphic functions have a decomposition 
    \begin{equation}
    \label{eq:holchaosdecomposition}
    \Psi_{Hol}(\phi^{\mathbf{x}})=\sum_{n=0}^\infty\psi^{n}_{\vec{x}_n}(\phi^n)^{\vec{x}_n}
    \end{equation}
     independently of the gaussian measure considered. 


    For the real case we also have a Wiener-Ito decomposition.  In this case we work out the decomposition form the real $S$ transform and therefore Wick monomials are given by 


    \begin{equation}
    \label{eq:wickMonomialsGenerator}
    :\varphi^n:\lvert_\Delta^{\vec{\mathbf{x}}}=
    \prod_{i=1}^n\frac{\delta}{\delta \xi_{y_i}}
    \exp{\big(\xi_x\varphi^x-\frac12\xi_x\Delta^{xy}\xi_y\big)} \lvert_{\xi_{\mathbf{x}}=0}
    \end{equation}
    or recursively 
\begin{definition}
    \label{def:wickMonomials}{\bf (Wick monomials)} 
    We denote $:\varphi^n:\lvert_\Delta^{\vec{\mathbf{x}}}=\ :\varphi^n:\lvert_\Delta^{\mathbf{x_1,\cdots,x_n}}$, and call Wick monomial of degree $n$ with respect to the covariance $\Delta^{\mathbf{xy}}$, the distribution belonging to the symmetric tensor product of $n$ copies on $\mathcal{N}'$ i.e. $:\varphi^n:\lvert_\Delta^{\vec{\mathbf{x}}}\in(\mathcal{N}')^{n\hat{\otimes}}$ such that
    \begin{align}
        :\varphi^0:\lvert_\Delta&=1\nonumber\\
        :\varphi^1:\lvert_\Delta^{\mathbf{x}}&=\varphi^{\mathbf{x}}\nonumber\\
        :\varphi^n:\lvert_\Delta^{\vec{\mathbf{x}},{\mathbf{x}}_{n-1},{\mathbf{x}}_n}&=
        :\varphi^{n-1}:\lvert_\Delta^{\vec{{\mathbf{x}}},{\mathbf{x}}_{n-1}}\varphi^{{\mathbf{x}}_n}-(n-1)\Delta^{\mathbf{x}_n(\mathbf{x}_{n-1}}:\varphi^{n-2}:\lvert_\Delta^{\vec{\mathbf{x}})}
        \label{eqn:wickPowers}
    \end{align}
    where the parenthesized indices are symmetrized according to the convention 
    \begin{equation}
    \label{eq:symmetrization}
    \varphi_{(x_1,\cdots,x_n)}:=\frac{1}{n!}\sum_{\pi\in S_n}\varphi_{x_{\pi(1)},\cdots,x_{\pi(n)}}
    \end{equation}
\end{definition}
Notice that this recursive definition coincides with the usual prescription for normal ordering in quantum field theory where $\Delta^{\mathbf{xy}}$ is the propagator of the theory considered.  Then the Wiener-Ito decomposition is 

$$\Psi(\varphi^\mathbf{x})=\sum_{n=0}^\infty \psi^{(n)}_{\vec{x_n}}:\varphi^n:\lvert_\Delta^{\vec{{x}_n}}$$
 the coefficients belong to the weighted Bosonic Fock space obtained from $\Gamma\mathcal{H}_\Delta$ with weights $\sqrt{n!}$ and the scalar product is 
\begin{multline}
    \label{eq:complexchaosprod}
    \int_{\mathcal{N}}D\mu(\varphi^\mathbf{x})\overline{\Phi}(\varphi^\mathbf{x})\Psi(\varphi^\mathbf{x})
    = \\   \sum_{n=0}^\infty n!\ \overline{\phi^{(n)}}_{\vec{u}_n}\big[\Delta^n\big]^{\vec{u}_n\vec{x}_n}  \psi^{(n)}_{\vec{x}_n}
    = 
    \sum_{n,}^\infty n! \big\langle\phi^{(n)}_{\vec{\mathbf{x}}_n},\psi^{(n)}_{\vec{\mathbf{x}}_n}\big\rangle_{\mathcal{H}^{n}_\Delta}
    \end{multline}
    \subsubsection{The Seagal-Bargmann tranform: Isomorphisms between the real and holomorphic cases.}
    In the case where the covariance for the complex case, $\Delta^{\mathbf{xy}}$ is real and symmetric we can stablish a unitary isomorphism between the subspace of holomorphic functions $L^2_{Hol}(\mathcal{N}'_{\mathbb{C}},D\mu_c)$ and  $L^2(\mathcal{N}',D\mu)$ just by performing an integration over the imaginary part of the coordinates:    


    \begin{definition}
        \label{def:Segaltranf}{\bf (Segal-Bargmann transform)} We call Segal-Bargmann transform the unitary isomorphism
        $$
        \begin{array}{cccc}
        \mathcal{B} :&  L^2(\mathcal{N}',D\mu)  & \longrightarrow &  L_{Hol}^2(\mathcal{N}_{\mathbb{C}}',D\mu_c) \\ 
         & \Psi_{Real} &\longmapsto & \Psi_{Hol}  
        \end{array}
        $$
        where both $\mu,\mu_c$ are gaussian measures with a real symmetric covariance $\Delta^{\mathbf{xy}}$, adapted in each case to its domain. To write the explicit formula for this isomorphism we follow a slightly different convention and write $\phi^\mathbf{x}=\varphi^\mathbf{x}+i\pi^\mathbf{x}$ to simplify computations. Then it follows   $\int_{\mathcal{N}'_\mathbb{C}} D\mu_c(\phi^{\mathbf{x}})=\int_{\mathcal{N}'} D\mu(\varphi^{\mathbf{x}})\int_{\mathcal{N}'} D\mu(\pi^{\mathbf{x}})$. Let $\Psi_{Hol}$ be a holomorphic function with chaos decomposition given by \eqref{eq:holchaosdecomposition}, integrating over the imaginary degrees of freedom 
    \begin{equation*}
        \label{eq:schHolrelation}
        \begin{split}
        \int D\mu(\pi^{\mathbf{x}})\Psi_{Hol}(\phi^\mathbf{x})=
        %\\ 
        \sum_{n=0}^\infty &\int D\mu(\pi^{\mathbf{x}}) \psi^{(n)}_{\vec{x}_n}{[(\varphi+%&
        i\pi)^n]^{\vec{x}_n}}=
        \\
        \sum_{n=0}^\infty  \psi^{(n)}_{\vec{x}_n}\frac{\partial}{\partial \xi^n_{\vec{x}_n}}\int D\mu(\pi^{\mathbf{x}})& e^{\xi_z(\varphi+i\pi)^z}\bigg\lvert_{\xi_{\mathbf{x}}=0}= 
        \\
        \sum_{n=0}^\infty  \psi^{(n)}_{\vec{x}_n}
        \frac{\partial}{\partial \xi^n_{\vec{x}_n}}\Bigg[&   e^{{\xi_z\varphi^z} } \int D\mu(\pi^{\mathbf{x}}) e^{i{\xi_z\pi^z}}\Bigg]_{\xi_{\mathbf{x}}=0}= 
        \\
        \hspace{ 2 em}\sum_{n=0}^\infty \psi^{(n)}_{\vec{x}_n}\frac{\partial^n}{\partial \xi^n_{\vec{x}_n}}e^{\xi_x\varphi^x-\frac{\xi_u\Delta^{uv}\xi_v}{2}}\Big\lvert_{\xi_\mathbf{x}=0}&=
        \sum_{n=0}^\infty \psi^{(n)}_{\vec{x}_n}:\varphi^n:\lvert_{{\Delta}}^{\vec{x}_n}=
         \Psi_{Real}(\varphi^{\mathbf{x}})
        \end{split}
    \end{equation*}
    where $\xi_\mathbf{x}\in \mathcal{N}$ is real.   With this consideration,  $\mathcal{B}(\Psi_{Real})$ is computed obtaining the chaos decomposition of  $\Psi_{Real}$ and replacing Wick monomials by regular holomorphic monomials. 
    
    A similar isomorphism is obtained using the antiholomorphic subspace in the analgous form
    
    $$
    \overline{\mathcal{B}} :L^2(\mathcal{N}',D\mu)   \longrightarrow  L_{\overline{Hol}}^2(\mathcal{N}_{\mathbb{C}}',D\mu_c).
    $$
\end{definition}





\subsubsection{Malliavin calculus: creation and annihilation operators and charges}

The space $L^2(\mathcal{N}',D\mu)$ is endowed with some useful properties that may be used to model very relevant tools in QFT such as creation and annihilation operators, or the number operator. Let us briefly introduce them in the (simple) case of real fields and later we will generalize it to the case of complex ones.

We start introducing  polynomials with coefficients in $\mathcal{N}$. A polynomial of  a variable $\varphi^{\mb{x}}\in \mathcal{N}'$ and degree $n$, which we denote as $P_n(\varphi^\mb{x})$,  can be written as
\begin{equation}
\label{eq:polynomials} 
P_n(\varphi)= C^0+C^1_x\varphi^x+C^{2}_{x_1x_2}(\varphi^2)^{x_1x_2}+\cdots+C^n_{\vec{x}_n}(\varphi^n)^{\vec{x_n}}
\end{equation}    

Where $C^0\in\mathbb{C}$ and the coefficients, for $n>0$, belong to the symmetric tensor products of $n$ copies of $\mathcal{N}$ i.e. $C^n_{\vec{\mathbf{x}}_n}\in \mathcal{N}^{n\hat{\otimes}}$. The sets of polynomials are
\begin{gather}
\label{eq:degnpolynomials}
\mathcal{P}_n(\mathcal{N}')=\left\{P(\varphi^{\mb{x}})\mid P\textrm{ polynomial of degree} \leq n\right\}\hspace{0.5 cm}
\mathcal{P}(\mathcal{N}')= \bigcup_{n=0}^\infty \mathcal{P}_n.
\end{gather}

We know that $\mathcal{P}(\mathcal{N}')$ is dense in $L^2(\mathcal{N}',D\mu)$ so we start introducing the operations of Malliavin calculus on $\mathcal{P}(\mathcal{N}')$ and then extend them to the whole space by linearity. On polynomials, the Frechèt derivative has an expression, particularly useful for practical purposes, in terms of the operator $\partial_{\varphi^\mathbf{y}}=\partial_{\mathbf{y}}$ that can be defined over monomials as 
\begin{equation}
\label{eq:Frechetderivative}
\partial_{\varphi^y} C_{x_1,\cdots,x_n}(\varphi^n)^{x_1,\cdots,x_n}=nC_{y,x_1,\cdots,x_{n-1}}(\varphi^{n-1})^{x_1,\cdots,x_{n-1}}
\end{equation}
and extended by linearity to $(\mathcal{N})$ without further complications.

The Malliavin derivative is the extension of the Frechèt derivative, defined in \eqref{eq:Frechetderivative}, to an operator 
\begin{equation}
\label{eq:malliavinDerivative}
\partial_\mathbf{x}: \mathbb{D}_\mu^{2,1}\longrightarrow \mathcal{H}_\Delta\otimes L^2(\mathcal{N}',D\mu)
\end{equation}
where the domain of the operator is
\begin{equation}
    \label{eq:MalliavinSovoleb}
    \mathbb{D}_\mu^{2,1}=\left\{\sum_{n=0}^\infty\psi^{(n)}_{\vec{x}_n}:\varphi^n:\lvert_\Delta^{\vec{x}_n}\textrm{ s.t } \sum_{n=1}^\infty n\cdot n!\lVert\psi^{(n)}_{\vec{\mathbf{x}}_n}\rVert_{\Delta}^2< \infty \right\}
\end{equation}
and acts over $\Psi(\varphi^\mathbf{x})$ as

\begin{equation}
\label{eq:actionMalliavin}
\partial_x\Psi(\varphi^\mathbf{x})=\sum_{n=1}^\infty n\psi^{(n)}_{x,\vec{y}_{n-1}}:\varphi^{n-1}:\lvert_\Delta^{\vec{y}_{n-1}}
\end{equation}
The key feature of this derivative is that the resulting derivative makes sense as a $\mathcal{H}_\Delta$ function and does not have a pointwise meaning as  \eqref{eq:Frechetderivative}. The Malliavin derivative can thus be considered as the generator of derivatives in directions of the Cameron-Martin Hilbert space. Furthermore, through the isomorphism $\mathcal{I}$, the contraction $\Delta^{xy}\partial_y$ is mapped to the annihilation operator of the Fock space $\Gamma(\mathcal{H}_\Delta)$.

The adjoint operator of the Malliavin derivative is the Skorokhod integral $\partial^{*\mathbf{x}}$. It is an operator defined over $\textrm{Dom}(\partial^{*\mathbf{x}})\subset\mathcal{H}_\Delta\otimes L^2(\mathcal{N}',D\mu)$ 

\begin{equation}
    \label{eq:SkorokhodIntegral}
    \begin{array}{cccc}
        \partial^{*\mathbf{x}}:& \textrm{Dom}(\partial^{*\mathbf{x}}) &\longrightarrow &L^2(\mathcal{N}',D\mu)\\
         & \Psi_{\mathbf{x}}(\varphi^{\mathbf{y}}) & \longmapsto & \partial^{*{x}} \Psi_{x}(\varphi^{\mathbf{y}})
    \end{array}
\end{equation}
It is important to remark that the Skorokhod integral is able to integrate more elements than those obtained from \eqref{eq:actionMalliavin}. Indeed, a general element $\Psi_{\mathbf{x}}\in\mathcal{H}_\Delta\otimes L^2(\mathcal{N}',D\mu)$ admits a Wiener-Ito decomposition in which 
\begin{equation}
\label{eq:parametricWienerIto}
\sum_{n=0}^\infty\psi^{(n)}_{x,\vec{y}_n}:\varphi^n:\lvert_\Delta^{\vec{y}_n}\hspace{0.25 cm} \textrm{ in general } \hspace{0.25 cm}  \psi^{(n)}_{(x,\vec{y}_n)}\neq \psi^{(n)}_{x,\vec{y}_n}
\end{equation}
In other words $(\psi^{(n)}_{x,\vec{y}_n})_{n=0}^\infty$ is not necessarily symmetrized and hence it may not be an  element of the Bosonic Fock space $\Gamma(\mathcal{H}_\Delta)$. Nonetheless if  
$\tilde{\psi}^{(n+1)}_{\vec{\mathbf{z}}_{n+1}}=\psi^{(n)}_{(\mathbf{x},\vec{\mathbf{y}}_n)}$ is indeed the decomposition of an element of $\Gamma(\mathcal{H}_\Delta)$ then $\Psi_{\mathbf{x}}\in  \textrm{Dom}(\partial^{*\mathbf{x}})$ and 

\begin{equation}
\label{eq:actionSkorokhod}
\partial^{*{x}}\Psi_{{x}}=\sum_{n=1}^{\infty} \tilde{\psi}^{(n)}_{\vec{z}_n}:\varphi^n:\lvert_\Delta^{\vec{z}_n}
\end{equation}
This, for a large class of functions in $\textrm{Dom}(\partial^{\star\mathbf{x}})$, can be written in terms of the Malliavin derivative as 
\begin{equation}
\label{eq:creationOperator}
\partial^{*{x}}=-\Delta^{xy}\partial_y+\varphi^x
\end{equation}

The Skorokhod integral is a generalization of the Ito integral, that coincides with it when $d=1$ and we consider only $\Psi_x$ as a process adapted to the filtration of Brownian motion, see \cite{Nunno2009}.  This operator is mapped onto the creation operator in the Fock space picture by the isomorphism $\mathcal{I}$. 
It is even possible to extend these operations to the space of Hida distributions  $(\mathcal{N}_\mathbb{C})'$ that are treated in detail in \cite{alonso2023a}. For further details on Hida-Malliavin calculus,  we refer the interested reader to  \cite{Nunno2009}.

Another operator that can be easily written in terms of the latter two is the particle number 

\begin{equation}
\label{eq:particleNumber}
N_{\Delta}=\partial^{*x}\partial_x=-\Delta^{xy}\partial_x\partial_y+\varphi^x\partial_x
\end{equation}

These notions can be easily extended to the holomorphic and antiholomorphic cases, which will represent matter-antimatter  models. 
\begin{gather}
\label{eq:holomorphicMalliavin}
\partial_{\phi^x}\Psi(\phi^{\mathbf{x}},\bar\phi^\mathbf{y})=\sum_{n,m=0}^\infty n\mathcal{W}_\Delta(\phi^{n-1}\bar{\phi}^m)^{\vec{z}_{n-1}\vec{y}_m}\psi^{(n,\bar{m})}_{x,\vec{z}_{n-1}\vec{y}_m} \\
\partial_{\bar\phi^x}\Psi(\phi^{\mathbf{x}},\bar\phi^\mathbf{y})=\sum_{n,m=0}^\infty m\mathcal{W}_\Delta(\phi^{n}\bar{\phi}^{m-1})^{\vec{z}_{n}\vec{y}_{m-1}}\psi^{(n,\bar{m})}_{\vec{z}_{n}\vec{y}_{m-1},x} 
\end{gather}
Here each derivative must be understood in the Wirtinger sense 

\begin{equation}
\label{eq:wirtinger}
\partial_{\phi^x}=\frac{1}{\sqrt{2}}\big(\partial_{\varphi^x}-i \partial_{\pi^x} \big),\hspace{0.5 cm} \partial_{\bar\phi^x}=\frac{1}{\sqrt{2}}\big(\partial_{\varphi^x}+i \partial_{\pi^x} \big)
\end{equation}
and also the Skorokhod integrals  

\begin{align}
\label{eq:holomorphicSkorokhod}
\partial^{*\phi^x} =-\Delta^{xy}\partial_{\bar\phi^y}+\phi^x, & &\partial^{*\phi^x}\mathcal{W}_\Delta(\phi^{n}\bar{\phi}^m)^{\vec{z}_{n}\vec{y}_m}= \mathcal{W}_\Delta(\phi^{n+1}\bar{\phi}^m)^{x,\vec{z}_{n}\vec{y}_m}\nonumber \\
\partial^{*\bar\phi^x} =-(\Delta^*)^{xy}\partial_{\phi^y}+\bar\phi^x, & &\partial^{*\bar\phi^x}\mathcal{W}_\Delta(\phi^{n}\bar{\phi}^m)^{\vec{z}_{n}\vec{y}_m}= \mathcal{W}_\Delta(\phi^{n}\bar{\phi}^{m+1})^{\vec{z}_{n}\vec{y}_m,x}
\end{align}

Then we can write down some operators, the number of particles $N$ and the $U(1)$ charge $Q$ encoding the matter/antimatter content of the states:
\begin{align}
\label{eq:ChargeAndNumber}
    N=\partial^{*\phi^x}\partial_{\phi^x}+ \partial^{*\bar\phi^x}\partial_{\bar\phi^x} & & Q=\partial^{*\phi^x}\partial_{\phi^x}- \partial^{*\bar\phi^x}\partial_{\bar\phi^x}
\end{align}

With those operators we see that the decomposition of \eqref{eq:complexchaosDecomposition} is such that

\begin{equation}
\label{eq:descompositionCharge}
L^2\big(D\mu_c,{\mathcal{N}_{\mathbb{C}}'}\big)=\sum_{n=0}^\infty\oplus \left( \sum_{m=0}^n\oplus \mathcal{H}_\Delta^{(n-m,\bar{m})} \right)
\end{equation}
where $\mathcal{H}_\Delta^{(n-m,\bar{m})}$ are simultaneous eigenspaces of $N$ and $Q$ with eigenvalues $n$ and $n-2m$ respectively. 


\section{Adapting Classical Field Theory to quantization}
\label{sec:CFTQuantizar}
Our first task is to adapt to quantization the geometrical description of the classical field theory that we are going to quantize. Alternative geometrical descriptions of classical field theories in terms of jet bundles also exist and admit  well defined Hamiltonian formalisms. Nonetheless, our goal is to build, from this classical manifold, a geometrical description of the space of quantum states on the space of fields, and our choice provides us with a simpler analog of the finite dimensional construction for quantum mechanical systems (see \cite{Kibble1979}).

We will focus on the real scalar field, i.e. the phase space of the theory will be given by the cotangent bundle to a real line bundle $\pi_{\Sigma L}:L\to \Sigma$. For simplicity we will consider the trivial case in which $L=\mathcal{N}$ is a space of smooth functions over $\Sigma$.

Furthermore, the set of fields modelled by any of these function spaces should be solutions of a certain dynamical field equation representing the type of field. The simplest example will be the case of a Klein-Gordon field theory. In this work we will consider only linear equations and therefore the corresponding manifold of fields will admit a linear structure which will significantly simplify the geometrical setting. 

\subsection{Choosing the space of classical fields and their momenta}

If $\mathcal{N}$ represents the set of fields, the corresponding phase space should correspond to the cotangent bundle $T^*\mathcal{N}\simeq \mathcal{N}\times \mathcal{N}'$, where we used the linearity of $\mathcal{N}$ and the corresponding triviality of the cotangent bundle. The coordinate representation of an element of this space is given by $(\varphi_\mathbf{x},\pi^\mathbf{y})$, with the coordinate systems described above. This space is endowed with a weakly symplectic form  which takes the local expression:

\begin{equation}
\label{eq:symplecticForm}
\omega_{T^*\mathcal{N}}=  d\pi_x \wedge d\varphi^x.
\end{equation}

This situation is problematic in the sense that the model space for the fiber and the base of the cotangent bundle $T^*\Gamma(\Sigma,L)$ are different. To quantize the system is convenient to express both coordinates over the same model space. To do so we endow $\Sigma$ with a Riemannian structure $h$, use it to define a measure $d\mathrm{Vol}$ and use the natural structure of rigged Hilbert space given by 
\begin{equation}
\label{eq:rhsOnSigma}
\mathcal{N}\subset L^2(\Sigma,d\textrm{Vol})\subset \mathcal{N}'
\end{equation} 

As a result, the elements of $\mathcal{N}$ corresponds to a linear subspace of its dual and we can write the coordinate representation of the image of an element of $\mathcal{N}$ as $\delta^{xy}\varphi_x$. Thus, in this setting, the manifold of fields must  be $\mathcal{M}_F=\mathcal{N}'\times \mathcal{N}'$ and coordinates $(\varphi^\mathbf{x},\pi^\mathbf{y})$ are to be understood as Darboux coordinates for a densely defined weakly symplectic structure
\begin{equation}
\label{eq:symplecticFormlifted}
\omega_{\mathcal{M}_F}=\delta_{xy} d\pi^x \wedge d\varphi^y =\int_\Sigma \frac{d^dx}{\sqrt{\lvert h\rvert}}\ [d\pi(x)\otimes d\varphi(x)-d\varphi(x)\otimes d\pi(x)],
\end{equation}
where now $\pi(x),\varphi(x)$ are densities. Notice that this symplectic structure is defined over the dense subspace $L^2(\Sigma,d\textrm{Vol})\subset \mathcal{N}'$. If we consider the space of polynomial functions  $\mathcal{P}(\varphi^\mathbf{x},\pi^\mathbf{y})$ with coefficients in $\mathcal{N}$ ( that, eventually, can be  extended to more general completions) we can also consider a Poisson bivector of the form:
\begin{multline}
\label{eq:poissonbracket}
\left\{P,Q\right\}_{\mathcal{M}_F}=
\delta^{xy}\left( \partial_{\varphi^x}P\partial_{\pi^y}Q-\partial_{\pi^x}P\partial_{\varphi^y}Q \right):=\\
\int_\Sigma d\textrm{Vol}(x)\left( \frac{\partial P}{\partial\varphi(x)}\frac{\partial Q}{\partial\pi(x)}-\frac{\partial P}{\partial\pi(x)}\frac{\partial Q}{\partial\varphi(x)} \right),
\end{multline}
where partial derivatives are understood as Frechèt derivatives, which are well defined for $\mathcal{P}(\varphi^\mathbf{x},\pi^\mathbf{y})$. This Poisson structure is often mistaken with a symplectic structure in QFT quantization schemes since, for linear spaces, the difference is subtle see \cite{Brunetti2015} chapter 5 for further discussion on that matter. With this Poisson bracket at hand we can proceed as usual and write down Hamiltonian dynamics of classical observables.  Classical observables are elements of a suitable completion of the set of polynomials in the field and momenta variables $\mathcal{P}(\varphi^\mathbf{x},\pi^\mathbf{y})$ and existence and uniqueness of such dynamics is, in general, strongly dependent on this completion. Such problems, though, will not be covered in this work.

Summarizing: our classical phase space shall be considered to be the linear subspace of $\mathcal{M}_F=\mathcal{N}'\times \mathcal{N}'$ corresponding to the injection of $\mathcal{N}$ by the rigged Hilbert space structure introduced on $L^2(\Sigma, d\textrm{Vol})$ by the riemannian structure of the Cauchy hypersurface $\Sigma$. On this space we can consider two (almost) equivalent structures: a (weakly) symplectic structure $\omega_{\mathcal{M}_F}$ and a Poisson tensor $\{ \cdot, \cdot \}_{\mathcal{M}_F}$.  

It is important to remark a subtle aspect of the construction: while the symplectic form on $T^*\mathcal{N}$ (Equation \eqref{eq:symplecticForm}) is canonical, the one on $\mathcal{M}_C$ (Equation \eqref{eq:symplecticFormlifted}   is not, since it depends on the choice of measure $d\mathrm{Vol}$ which we used to define the Gel'fand triple. The same is true for the Poisson bracket \eqref{eq:poissonbracket}. This aspect will be important when considering quantum classical hybrid models.


Now, we will learn to introduce another structure on it to define more conveniently the quantization of the classical model. In order to do that, we need to introduce a complex structure on $\mathcal{M}_F$.

\subsection{Introducing a complex structure}
In order to obtain a system better  suited for geometric quantization we need an additional (almost) complex structure $J_{\mathcal{M}_F}$ satisfying:
\begin{itemize}
	\item it is densely defined,
	\item  it leaves $D(J_{\mathcal{M}_F})\subseteq L^2(\Sigma,d\textrm{Vol})\subset \mathcal{N}'$ invariant
	\item and such that $J_{\mathcal{M}_F}^2=-\mathds{1}$ in its domain $D(J_{\mathcal{M}_F})$.
\end{itemize}

We shall see below that the complex structure required for quantization will be determined by the classical Hamiltonian on the space of fields. Hence, in order to consider different possible classical dynamics, we must explore the space of possible complex structures on $\mathcal{M}_F$.  It is simple to verify that there are infinitely many different structures satisfying those conditions. To express the most general complex structure let $(\varphi^{\mathbf{x}},\pi^{\mathbf{y}})$ be Darboux coordinates for $\mathcal{M}_F$ and let us consider the matrix elements $(A^t)^{\mathbf{x}}_\mathbf{y}:= \delta_{u \mathbf{y}} A^u_v\delta^{v\mathbf{x}}$. Notice that the definition of these variables already depends on the choice of the Hilbert space structure for the Gel'fand triple.

An arbitrary (almost) complex structure on $\mathcal{M}_F$ is locally expressed as 
\begin{multline}
\label{eq:complexStructure}
-J_{\mathcal{M}_F} =
(\partial_{\varphi^y},\partial_{\pi^y})\left( \begin{array}{cc}
A^y_x & \Delta^y_x \\
D^y_x & -(A^t)^y_x
\end{array} \right)  \left( \begin{array}{c}
d\varphi^{x} \\
d\pi^{x}
\end{array} \right)=\\
=  \partial_{\varphi^y} \otimes[A^y_xd\varphi^x+\Delta^y_xd\pi^x]+ \partial_{\pi^x}\otimes [D^y_xd\varphi^x-(A^t)^y_xd\pi^y]
\end{multline}
With $(A\Delta)^x_y=A^x_z\Delta^z_y$ we have $A^2+\Delta D=-\mathds{1},\ \Delta^t=\Delta, D^t=D, A\Delta=\Delta A^t$ and $A^tD=DA$. Notice that $D$ is fixed once $\Delta$  and $A$ are known. Let $K^{\mathbf{x}}_{\mathbf{y}}$ be the inverse of $\Delta_\mathbf{y}^{\mathbf{x}}$, i.e. $\Delta^{x}_zK^z_{y}=\delta^x_y$ then $D=(iA^t+\mathds{1})K(iA-\mathds{1})$.

 

This complex structure, together with the symplectic form \eqref{eq:symplecticFormlifted} induces a (pseudo)-Riemannian structure $\mu_{J,\mathcal{M}_F}(\cdot,\cdot)=\omega_{\mathcal{M}_F}(\cdot,-J\cdot)$, which, in the coordinates above, reads: 
\begin{multline}
\label{eq:riemannianInduced}
\mu_{\mathcal{M}_F}=(d\varphi^{x},d\pi^{x}) \left( \begin{array}{cc}
    -D_{xy} & A^t_{xy} \\
    A_{xy} & \Delta_{xy}
    \end{array} \right) \left( \begin{array}{c}
        d\varphi^{y} \\
        d\pi^{y}
    \end{array} \right)=\\
    -D_{xy}d\varphi^x\otimes d\varphi^y+\Delta_{xy}d\pi^x\otimes d\pi^y+ A_{xy}(d\varphi^y\otimes d\pi^x+d\pi^x\otimes d\varphi^y)
\end{multline}
where $A_{xy}=\delta_{xz}A^z_y$. To obtain a Riemannian structure form here we must impose also $\Delta_{\mathbf{xy}}>0>D_{\mathbf{xy}}$. 

With a change of coordinates, locally given by
\begin{align}
\label{eq:changeOfCoordinates}
d{\tilde{\varphi}^x}&=d{\varphi^x} &
d{\tilde{\pi}^x}&=A^x_yd{\varphi^y}+\Delta^x_yd{\pi^y} 
\\
\partial_{\tilde{\varphi}^x}&=\partial_{\varphi^x}- (K A)^y_x\partial_{\pi^y} &
\partial_{\tilde{\pi}^x}&=K^y_x\partial_{{\pi}^y}
\end{align}
 we can put in a canonical way $J_{\mathcal{M}_F}=\partial_{\tilde{\pi}^x}\otimes d\tilde{\varphi}^x-\partial_{\tilde{\varphi}^x}\otimes d\tilde{\pi}^x$, $\mu_{\mathcal{M}_F}=K_{xy}(d\tilde{\varphi}^x\otimes d\tilde{\varphi}^y+d\tilde{\pi}^x\otimes d\tilde{\pi}^y)$ and $\omega_{\mathcal{M}_F}=K_{xy} d\tilde{\pi}^x\wedge d\tilde{\varphi}^y$. Notice that we can cast the structures into this canonical form with other changes of variables, but we choose this one because it leaves $\varphi^\mathbf{x}$ (the manifold of fields) invariant. Notice that we are transforming the original cotangent bundle structure we began with, using its triviality, as the new coordinates combine base and fiber coordinates of the original bundle structure.

Having chosen a complex structure, we can define an isomorphism for $\mathcal{M}_F$ to become a complex vector space $\mathcal{M}_C=\mathcal{M}_F^{\mathbb{C}}$. Notice that while $\mathcal{M}_F$ depends on the choice of the Hilbert space structure for the Gel'fand triple (and the original $T^*\mathcal{N}$), the definition of $\mathcal{M}_C$ depends also on the choice of the complex structure.  This fact will be important later. Notice, anyway, that the symplectic structures on $T^*\mathcal{N}'$, $\mathcal{M}_F$ and $\mathcal{M}_C$ are all diffeomorphic, but the diffeomorphisms depend on the metric structure $h$ on $\Sigma$ (for $\mathcal{M}_F$)  and the complex structure $J_{\mathcal{M}_C}$ (for $\mathcal{M}_C)$. 

 Having chosen the complex manifold $\mathcal{M}_C$, we choose holomorphic and antiholomorphic  coordinates which are locally written as $(\phi^{\mathbf{x}},\bar{\phi}^{\mathbf{y}})=\frac{1}{\sqrt{2}}(\tilde{\varphi}^{\mathbf{x}}+i\tilde{\pi}^{\mathbf{x}},\tilde{\varphi}^{\mathbf{y}}-i\tilde{\pi}^{\mathbf{y}})$. In this new manifold the complex structure becomes  $J_{\mathcal{M}_C}$, which  is now written as: 

\begin{equation}
\label{eq:complexholom}
J_{\mathcal{M}_C}=i \partial_{\phi^x}\otimes d\phi^x-i \partial_{\bar{\phi}^x}\otimes d\bar{\phi}^x.
\end{equation}
At the tangent space level we can always write the projection operators $P_{\pm}=\frac12(\mathds{1}\mp iJ)$ such that $P_{+}\partial_{\phi^{\mathbf{x}}}=\partial_{\phi^{\mathbf{x}}}$ and $P_{-}\partial_{\bar\phi^{\mathbf{x}}}=\partial_{\bar\phi^{\mathbf{x}}}$. We shall use these projections in the following sections.

As the triple $(\mu_C, \omega_C, J_C)$ defines a Kähler structure on $\mathcal{M}_C$, there exists
a (non unique) globally defined on a dense subspace (because the manifold is linear) Kähler potential  $\mathcal{K}$

\begin{equation}
\label{eq:kahlerpotential}
\mathcal{K}(\bar \phi^x, \phi^x)=\bar \phi^x  K_{xy} \phi^y .
\end{equation}



From the potential we can define the corresponding  Kähler form:

\begin{equation}
\label{eq:kahlerform}
\omega_{\mathcal{M}_C}= i\partial \bar \partial \mathcal{K},
\end{equation}
where $\partial$ and $\bar \partial$ represent the Dolbeault operators. It is important to notice that, as usually happens in the literature\footnote{See for example \cite{Woodhouse1997}}, the convention on the sign change from real to holomorphic coordinates, this is $\omega_{\mathcal{M}_C}=-\omega_{\mathcal{M}_F}$. With the Kähler form and the expression of $J_{\mathcal{M}_C}$ of \eqref{eq:complexholom} we obtain the expression of the riemannian form $\mu_{\mathcal{M}_C}$ and the hermitian form: $$h_{\mathcal{M}_C}= \frac{\mu_{\mathcal{M}_C}+i\omega_{\mathcal{M}_C}}{2}= K_{xy}d\phi^x\otimes d\bar{\phi}^y,$$ 
with inverse $h^{-1}_{\mathcal{M}_C}=\Delta^{xy}\partial_{\phi^x}\otimes\partial_{\bar\phi^y}$.

Finally, the expression of the Poisson bracket \eqref{eq:poissonbracket} in holomorphic coordinates becomes:
\begin{equation}
\label{eq:poissonHol}
\left\{P,Q\right\}_{\mathcal{M}_C}=
i\Delta^{xy}\left( \partial_{\phi^x}P\partial_{\bar\phi^y}Q-\partial_{\bar\phi^x}P\partial_{\phi^y}Q \right).
\end{equation}


It is important to notice that the election of the rigged Hilbert space made in \eqref{eq:rhsOnSigma}, to express the system in Darboux coordinates, is different for holomorphic coordinates and the sign of \eqref{eq:poissonbracket} is reversed due to our conventions. In the first case, we built the triple from the real classical fields onto the real momentum fields. In holomorphic coordinates the natural triple is defined from the complexified classical fields 
$\mathcal{N}_{\mathbb{C}}$ onto the complexified momentum fields $\mathcal{N}_{\mathbb{C}}'$:

\begin{equation}
\label{eq:HolomorphicTriple}
\mathcal{N}_{\mathbb{C}}\subset \mathcal{H}_\Delta\cong \mathcal{H}_K^*\subset {\mathcal{N}_{\mathbb{C}}'}
\end{equation}
 where $\mathcal{H}_\Delta=\overline{(\mathcal{N}_{\mathbb{C}},\langle \cdot,\cdot\rangle_\Delta)}$ is the  Hilbert space obtained by completion of $\mathcal{N}_{\mathbb{C}}$ with the scalar product induced by  the tensor $\Delta$ as $\langle \chi_{\mathbf{x}},\xi_{\mathbf{y}}\rangle_\Delta=\bar{\chi}_x\Delta^{xy}\xi_y$.  This space is identified with the dense subset of $\mathcal{N}_{\mathbb{C}}'$ on which the hermitian form $h_{\mathcal{M}_C}$ is finite. This Hilbert space  has a physical interpretation for $\mathcal{H}_\Delta$ as the space corresponding to the \textit{one-particle-state} structure. The space comprises the initial conditions for all the solutions of 'positive frequency' that are to be interpreted as particles in contrast with the elements of the conjugate $\mathcal{H}_{\Delta}^*$ which represent the corresponding anti-particles.    

On the other hand, we can consider a similar construction from the dual perspective if we consider the completion of  the dense subset of $\mathcal{N}_{\mathbb{C}}'$ on which the scalar product $\langle \cdot,\cdot\rangle_K$ induced by the hermitian tensor $K$ is defined.  Both spaces $\mathcal{H}_\Delta$ and $\mathcal{H}_K$ are identified by means of the Riesz representation theorem,  realizing the duality  $(\mathcal{N}_{\mathbb{C}},\mathcal{N}_{\mathbb{C}}')$ with the scalar product  $\langle \cdot,\cdot\rangle_\Delta$.  In the following, we will consider them as the same space.

Furthermore, the elements of $\mathcal{H}_K$ will be important later when we build the second quantization of our field theory. Indeed, when considering "wave-functions" on our manifold of classical fields, $\mathcal{H}_K$ is identified with the Cameron-Martin space represent the allowed translations of the unique Gaussian measure which has covariance $\Delta$. From that point of view, we will consider them to be the natural representation of the tangent space to our space of fields.

Summarizing: the classical phase space of fields is described as a Kähler complex manifold $\mathcal{M}_C$, which has associated a rigged Hilbert space structure encoding the physical \textit{one-particle-space} structure. In the next section, we will use these structures to define, by geometric quantization, a geometrical description of the quantum field theory corresponding to this  classical phase space of fields. 


\section{ Geometric quantization I: the general framework}
\label{sec:QFTCauchy}

The goal of this section is to build the quantum field theory corresponding to the classical model presented above. As our goal is a geometric description of the resulting quantum model, we will consider a geometric quantization framework adapted to our infinite dimensional manifold. The procedure is, roughly speaking, similar to the usual construction on finite dimensional symplectic manifolds (see \cite{Woodhouse1997}), and based in two steps.
\begin{itemize}
	\item a pre-quantization which builds an initial Hilbert space as the set of sections of a complex line bundle defined on the classical phase space and a representation of classical functions as operators on that set, 
	\item and the definition of a polarization, which chooses a Lagrangian submanifold on the phase space of fields, defines the Hilbert space as the space of sections restricted to the submanifold and the operators as those preserving it in a suitable way.
\end{itemize}
The first step is straightforward and relatively simple to define for a general symplectic classical manifold. The definition of the polarization is a more subtle task and introduces arbitrary choices which restrict the set of classical magnitudes that can be quantized.   We will present now, first the general construction adapted to our manifold of fields. Then, we will exploit the Kähler structure we defined on $\mathcal{M}_C$ to consider one particular type of polarization, the holomorphic polarization, which exhibits several nice properties and largely simplifies the technicalities of the construction. Finally, we will compare the resulting construction with the usual approach based on a real polarization.

\subsection{Prequantization I: The definition of the bundle and the measure}

The first step towards quantization is the complex prequantum line bundle $B$ over the classical field phase space. This is an hermitian line bundle $\pi_{\mathcal{M}_C,B}:B\rightarrow \mathcal{M}_C$, associated with a $U(1)$--principal bundle on $M$. On this bundle we define a principal connection which is required to have, as local curvature form, the symplectic form of $\mathcal{M}_C$. This can be done if we consider the local connection one form to be defined by the corresponding symplectic potential $\Theta$ (as our field manifold is a linear space, this choice is well defined globally). By using this connection, we can define a covariant derivative on the associated bundle $B$.


The next step is to find a prequantized Hilbert space $H$. To define this space we consider the set of square integrable sections $B$ with respect to a measure $D\mu$ invariant under symplectic transformations. Let us consider for simplicity the case where $\mathcal{M}_c$ can be covered with a chart in $\mathcal{N}'_\mathbb{C}$. In this case, as it is explained above, we will consider the unique Borel probability measure $D\mu_c(\phi^\mathbf{x})$ on $\mathcal{M}_C$ such that the characteristic functional $C$ obtained as  
\begin{equation}
\label{eq:gaussianmeasure}
C(\rho_{\mathbf{x}}, \bar \rho_{\mathbf{x}})=\int_{\mathcal{M}_C}D\mu_c(\phi^\mathbf{x} )e^{i (\bar\rho_x\phi^x+\rho_x\bar \phi^x) },
\end{equation}
satisfies:
\begin{itemize}
    \item $C(0,0)=1$
    \item $C$ is positive definite,
    \item and it is continuous in the Frechèt topology.
\end{itemize}

 
   With our choices of complex structure (and hence of Kähler potential and Cameron-Martin structure for $\mathcal{H}_\Delta$),  we can select the measure whose functional $C$ is equal to
\begin{equation}
\label{eq:C}
C(\rho_\mathbf{x}, \bar \rho_\mathbf{x})=e^{-h^{-1}_{\mathcal{M}_c}(\bar \rho_\mathbf{x}, \rho_\mathbf{x})}= e^{-\bar \rho_{x}\Delta^{xy} \rho_{y}}.
\end{equation}

From Minlos theorem, we know that there is one measure for each choice of a Frechèt-continuous functional for the exponent in the expression above.  We choose it to correspond to a function of the hermitian form   $h^{-1}_{\mathcal{M}_c}$, acting on  $\bar\rho_x d\phi^x$ as covectors, for the measure to be invariant under symplectic transformations. 

This measure can be explicitly constructed form the Gel'fand triple \eqref{eq:HolomorphicTriple}. We start with finite dimensional orthogonal discretizations of $\mathcal{H}_\Delta$. In this vision we perform a projective limit of the chain of orthogonal projections  $p_{nm}:\mathbb{C}   ^n\to \mathbb{C}^m,\ m<n$, in each of them we can define gaussian measures $d\mu_n$ expressed in the usual form with respect to the canonical Lebesgue measure.  If we perform the projective limit in the category of measure spaces we do not recover $\mathcal{H}_\Delta$ as a measure space, but we can instead perform the limit to obtain the Borel $\sigma$-algebra of  $\mathcal{N}'_\mathbb{C}$ and define the measure given by $C$ on it  (see \cite{Oeckl2012a,Gellfand} for further details). Due to this reason, we will call this measure \textit{gaussian}.


With this choice, the measure exists and then the scalar product of 
$\Psi^s,\Phi^s\in \Gamma(\mathcal{M}_C,B)$, sections of the prequantum bundle, can be written as:
\begin{equation}
\label{eq:scalarproductprequantum}
\langle\Phi^s,\Psi^s \rangle=\int_{\mathcal{M}_C}D\mu_c(\phi)\langle \Phi^s(\phi),\Psi^s(\phi)\rangle. 
\end{equation}

\subsection{Prequantization II. The states: introducing the vacuum}

\subsubsection{A trivializing section}
In this way we have introduced a measure on the set of sections of the prequantum bundle and we can define our quantum pure states to be the square-integrable sections with respect to the measure $D\mu_c$, or equivalently, those sections with finite norm. Now, we want to trivialize the bundle and factorize those sections with respect to a preferred one $\Psi_0$, which physically represents the vacuum state of the physical theory. We will ask this section to be constant in norm with respect to this measure, i.e.,
\begin{equation}
\label{eq:vacuum_norm}
\langle \Psi_0( \bar \phi^x, \phi^x), \Psi_0( \bar \phi^x, \phi^x) \rangle=1, \qquad \forall ( \bar \phi^x, \phi^x)\in \mathcal{M}_C.
\end{equation}

All other sections can then be represented as simple functions, which will be square integrable with respect to a measure defined by using $\Psi_0$. As sections are also functions on $\mathcal{M}_C$, we can write our pre-quantum Hilbert space as
\begin{equation}
\label{eq:preHilbert}
\mathcal{H}=\left \{ \Phi: \mathcal{M}_C, \to \mathbb{C} \text{ s. t. } \Phi\in L^2(\mathcal{M}_C,D\mu_c) \text{ as function } \right \}
\end{equation}


The local expression of the vacuum section can be further determined by the covariant derivative of the bundle. We can define the vacuum section $\Psi_0$ adapted to the symplectic potential (hence, to the complex structure) requiring that
\begin{equation}
\label{eq:vacuum}
\nabla_{X}\Psi_0=-i\theta(X)\Psi_0.
\end{equation}
Notice that while condition \eqref{eq:vacuum_norm} fixes the modulus of the vacuum state, this condition fixes the phase.  We shall get back to this point below, when considering the two types of polarization.  In any case, we can already remark that clearly the local expression of the state (or equivalently, the connection) depends on the choice of the complex structure on $\mathcal{M}_C$. 
Notice also that thinking in terms of the projective limit of discretized gaussian measures, $D\mu_c$ includes in its definition the expression of the term $\vert \Psi_0\vert^2$  which usually appears in the expressions in finite dimensions multiplying the Lebesgue measure. This is the reason for the necessity of  Equation \eqref{eq:vacuum_norm}, which only leaves free the vacuum state phase for the chosen measure.



Using this section $\Psi_0$ as a reference, any other section  $\Phi_s$ is given by $\Phi^s=\Phi\Psi_0$ where $\Phi:\mathcal{M}_C\to \mathbb{C}$ is a regular function of $L^2(\mathcal{N}'_\mathbb{C},D\mu_c)$. The action of the covariant derivative is then translated to the function $\Phi$ as:
\begin{equation}
\label{eq:covariantDerivative}
\nabla_X(\Phi\Psi_0)=\left[ X(\Phi)-i\theta(X)\Phi\right]\Psi_0.
\end{equation}



Our pre-quantum bundle is then identified with a space of regular functions of $L^2(\mathcal{N}'_\mathbb{C},D\mu_c)$ referred to a vacuum $\Psi_0$ rather than sections. The effect of the reference to the vacuum is that covariant derivatives incorporate an  extra multiplicative factor $-i\theta(X)$ provided by the symplectic potential to the directional derivative.   

Another key factor of this construction is that in the quantum theory every magnitude is obtained by integration of the classical degrees of freedom. This turns the classical manifold $\mathcal{M}_C$ into a rather auxiliary object. Let $\Psi$ be a pre-quantum state, and consider a translation by $\Psi(\phi^\mathbf{x})\mapsto \Psi(\phi^\mathbf{x}+h^\mathbf{x})$. Only when $h^\mathbf{x}\in \mathcal{H}_K\simeq \mathcal{H}_\Delta$, the Cameron-Martin Hilbert space, we can reabsorb this into the Gaussian measure \eqref{eq:gaussianmeasure} with an integrable Radon-Nikodym factor. Hence, we can consider $\mathcal{H}_\Delta$ as the only well-defined tangent vectors on the linear manifold $\mathcal{M}_C$ for the quantum counterpart, since they are the only ones that leave the Hilbert space invariant. Nonetheless, in general curved spacetimes this is not the case, and unitary evolution will not preserve  $L^2(\mathcal{N}'_\mathbb{C},D\mu_c)$ unless the spacetime is static or stationary.  We will deal with these issues in \cite{alonso2023a}.


\subsubsection{Ambiguity of the construction}
\label{sec:Ambiguity}

We may consider the existence of another connection on $B$ and the corresponding vacuum section $\tilde \Psi_0$ it is related to. As we want the section to be square integrable, we can consider the existence of a function $\alpha\in C^\infty(\mathcal{M}_C)$ satisfying

\begin{equation}
\label{eq:newvacuum}
\tilde \Psi_0=\alpha \Psi_0,
\end{equation}
for $\alpha:\mathcal{M}_C\to \mathbb{C}$ a function. Hence, sections of $B$ should satisfy
\begin{equation}
\label{eq:change-alpha}
\Phi_s= 
\begin{cases} 
    \tilde \Phi \tilde \Psi_0=\tilde \Phi \alpha\Psi_0 \\
    \Phi \Psi_0
\end{cases}
\end{equation}

For the two covariant derivatives to coincide on a section of $B$, we obtain that
\begin{equation}
\label{eq:vacua}
\nabla_X \tilde \Psi_0=
\begin{cases} 
  i\tilde \theta (X) \tilde \Psi_0 = i\tilde \theta (X)  \alpha \Psi_0  \\ 
  X(\alpha)\Psi_0+ \alpha(-i\theta(X) \Psi_0)
\end{cases}
\end{equation}
As this should happen for any vector field $X$, we conclude that the function $\alpha$ must satisfy
\begin{equation}
\label{eq:alphacondition}
d\alpha-i( \tilde \theta-\theta)\alpha=0
\end{equation}

Notice that as the modulus of the vacuum section is fixed to 1,  the function must be a phase $\alpha=e^{-if}$. This implies that the measure does not change but the symplectic potentials are required to satisfy the gauge condition
\begin{equation}
\label{eq:gauge}
\theta'-\theta=df.
\end{equation}
This is the freedom left from Equation \eqref{eq:vacuum_norm} once we fixed the measure $D\mu$. 
Notice that if the modulus of the function is equal to one the measure is not changed, since we are just choosing among the different phases of the vacuum state.



\subsection{Prequantization III: The quantization of observables}
The covariant differentiation is precisely the basic tool to build a quantization mapping $\mathcal{Q}$ for the classical functions $f\in C^\infty(\mathcal{M}_C)$ acting on the set of sections of $B$. Thus, the operator $\mathcal{Q}(f)$  becomes
\begin{equation}
\label{eq:quantization}
\mathcal{Q}(f)=-i\hbar \nabla_{X_f} +f,
\end{equation}
where $X_f$ is the Hamiltonian vector field associated to $f$. For the sake of simplicity we will take $\hbar=1$ in the rest of the paper.

When we consider the section $\Psi_0$, and the corresponding function $\Phi$, we can adapt $\mathcal{Q}$ to the new setting. Thus,  $\nabla_{X_f}$ must become a self-adjoint operator on the Hilbert space of square integrable functions, where  $\Phi$ is contained. 

The classical observables are chosen as functions on the classical phase space $F:\mathcal{M}_C \to \mathbb{R}$. 
\begin{equation}
\label{eq:Hamiltonina}
X_{F}\lrcorner \omega =-dF.
\end{equation} 
Remember that, being weakly symplectic, the equation above may not have a solution in general. 
But we will consider also that the first Malliavin derivative is defined for them and hence that $F\in\mathbb{D}_{\mu}^{2,1}$.   Therefore the expression
$$dF= \partial_{\phi^{x}}Fd\phi^x+\partial_{\bar\phi^{x}}Fd\bar\phi^x,$$
ensures that $dF\in \mathcal{H}_\Delta\otimes L^2(\mathcal{N}',D\mu) $ and $\mathcal{H}_\Delta$ is the natural domain of $\omega$ so we can always find $X_F\in \mathcal{H}_\Delta\otimes L^2(\mathcal{N}',D\mu)$. 
Then, its prequantum counterpart is an operator $\hat{F}:=\mathcal{Q}(F):H\to H$ given by 
\begin{equation}
\label{eq:actionPrecuantumOperator}
\mathcal{Q}(F) \Phi= F\Phi-i\nabla_{X_{F}}\Phi, \qquad \Phi\in H.
\end{equation}


It follows that this (pre)quantization procedure $\mathcal{Q}$  meets the Dirac quantization (pre)conditions \cite{hall2013}
\begin{enumerate}[label=Q\arabic*)]
    \item $\mathcal{Q}(F+G)=\mathcal{Q}({F})+\mathcal{Q}({G})$ 
    \item $\mathcal{Q}({F})\mathcal{Q}({G})-\mathcal{Q}({G})\mathcal{Q}({F})=-i\mathcal{Q}(\{F,G\}_{\mathcal{M}_F})$
    \item $\mathcal{Q}(F)=F\mathds{1}$ if $F$ is constant
\end{enumerate}
for any $F, G\in C^\infty(\mathcal{M}_C)$.

\subsection{The polarization}

The prequantization provides a Hilbert space too big to realize the quantum states since, in principle, $\Phi$ are functions with domain on the whole manifold while, in regular quantum theories, the domain must have half the degrees of freedom of the latter. 

To deal with this discrepancy we must restrict the domain of the functions $\Phi$ to a Lagrangian submanifold of $\mathcal{M}_{C}$. Notice, though, that in this setting the concept of Lagrangian submanifold is not as straightforward as in the finite-dimensional  case. First, we must remember that $\omega_{\mathcal{M}_C}$ is defined, at most, on a dense subset of the vector fields of  $\mathcal{M}_{C}$ and therefore that the condition of (co)isotropy must take into account this fact. 
 To bypass this difficulty we will introduce a polarization at the level of the \textit{one particle state structure} introduced in the complex description of the classical phase space \eqref{eq:HolomorphicTriple}. Thus, we will choose a Lagrangian subspace $P$ of the complex Hilbert space $P\subset H_{\Delta}$ on which the  symplectic form $\omega_{\mathcal{M}_C}$ is well defined. The quantum Hilbert space is given by:
 
 \begin{equation}
 \label{eq:ponarizedHilbert}
 \mathcal{H}=\overline{\left\{ \Psi^s \in H \textrm{ s.t. } \nabla_{\bar{X}}\Psi^s =0 \ \  \forall\  X\in P\right\}}
 \end{equation}

As, again, the space is linear, we can choose a symplectic potential $-d\Theta=\omega_{\mathcal{M}_C}$ adapted to the polarization such that 

 \begin{equation}
 \label{eq:adaptedPolar}
 \Theta(\bar{X})=0\ \forall \ X\in P
 \end{equation}

Under this assumption the definition of the polarization simplifies enormously since, due to \eqref{eq:covariantDerivative} we can identify the polarized Hilbert space as:

\begin{equation}
\label{eq:quantumFunctions}
\mathcal{H}=\overline{\left\{ \Psi \in C^\infty (\mathcal{M}_C) \textrm{ s.t. } d\Psi(\bar{X}) =0 \ \  \forall\  X\in P\right\}}
\end{equation}


Unitary dynamics is trickier to implement due to the fact that the quantization rule \eqref{eq:actionPrecuantumOperator} may produce operators that do not respect the polarization. Moreover the situation is even more complicated if we ask the representation of the quantum operators to be irreducible adding a fourth condition to Q1-Q3 above:

\begin{enumerate}[label=Q4)]
    \item {\bf (Irreducibility condition)} For a given set of classical observables $\{f_i\}_{\mathcal{I}}$ such that $\{f_i,g\}=0\ \forall i\in \mathcal{I}$ implies $g$ constant then if an operator $A$ commutes with every $Q(f_i)$ then $A$ is a multiple of the identity.
\end{enumerate}

Then it is known that no quantization satisfying Q1-Q4 exists. We can partially escape from those two problems choosing to quantize only a subset of classical observables such as the linear ones in Darboux coordinates $(f_x\varphi^x+g_x\pi^x)$. More complicated functions  exhibit the usual ordering problem of quantization.  We will deal with this issue in \autoref{sec:orderings}. 

\section{Geometric quantization II:  types of quantization}
\label{subsec:holom}  
Roughly speaking, we can consider two kind of polarizations: real polarizations, where the linear space $P$ is characterized by $P=P^*$ (which define what we call Schrödinger quantizations) and holomorphic ones for which $P\cap P^* =0$ (which define what we call holomorphic quantizations). Intermediate situations are combinations of those two extreme cases, at least in the finite dimensional setting. As the holomorphic case is easier to handle in many ways so we will begin our discussion with it. 

\subsection{Holomorphic quantization}

\subsubsection{Choosing holomorphic tensors}

For the holomorphic case the complex structure $J_{\mathcal{M}_C}$ (given by Equation \eqref{eq:complexStructure}) provides a projection $X\to \frac12(\mathds{1}-iJ_{\mathcal{M}_C})X$ for the elements of the Cameron Martin tangent space $X\in H_{\Delta}$, i.e., those elements in the tangent space which define integrable translations for the Gaussian measure $D\mu_c$. As $\mathcal{M}_C$ is a linear space,  the polarization is taken to be $P_{Hol}=\frac12(\mathds{1}-iJ_{\mathcal{M}_C}) H_{\Delta}$. The adapted symplectic potential in holomorphic coordinates $(\phi^{\mathbf{x}},\bar{\phi}^{\mathbf{y}})$ is obtained from the (non unique) Kähler potential $\mathcal{K}=\bar{\phi}^xK_{xy}{\phi}^y$ 
\begin{equation}
\label{eq:SymplecticPotential}
\Theta=-{i}\partial_{\phi^x} \mathcal{K}d\phi^x.
\end{equation}

The set of states of the quantum Hilbert space given by \eqref{eq:quantumFunctions} is therefore the space of holomorphic (square integrable) functions $\Phi:\mathcal{M}_C\sim \mathcal{N}_{\mathbb{C}}'\to \mathbb{C}$. These functions represent the excitation with respect to the vacuum state $\Psi_0$ (Equation \ref{eq:vacuum}), in such a way that the physical state corresponds to the section of the complex line bundle $\Phi\Psi_0$, which are square integrable with respect to the measure $D\mu_c$, or, equivalently, to the set of square integrable holomorphic functions if we take the section $\Psi_0$ as a reference:
\begin{equation}
    \label{eq:states_holomorphic}
    \mathcal{H}_{Hol}=L^2_{Hol}(\mathcal{N}_{\mathbb{C}}', D\mu_c).
    \end{equation}
This is, by construction, a well-defined Hilbert space but it is not the most frequent quantum model in QFT. Instead, it is more common to represent the states as element of a bosonic (or fermionic) Fock space.  We will discuss this point later.

\subsubsection{Segal isomorphism: defining Fock states}

From the properties analyzed above, we know that the Segal isomorphism presented in \autoref{def:SegalIso} allows us to identify the vectors in $L^2(\mathcal{N}_{\mathbb{C}}', D\mu_c)$ with the vectors of the Fock space constructed on the space of one-particle states $\mathcal{H}_\Delta$. As the quantum states correspond to holomorphic functions $\Psi\in L_{Hol}^2(\mathcal{N}_{\mathbb{C}}', D\mu_c)$, we can write:
\begin{equation}
\label{eq:Segal_holo}
\mathcal{I}(\Psi)=\left (\sqrt{n!}\psi^{(n,0)}_{\vec x_n} \right )_{n=0}^\infty, 
\end{equation}
where $\psi^{(n,0)}_{\vec x_n } $ correspond to the coefficients of the state $\Psi$ with respect to the space of Wick complex monomials

    \begin{equation}
        \label{eq:wienerItoSeagalwithTtransform2}
        \psi^{(n,0)}_{\vec{x}_n}=\frac{1}{n!}\prod_{i=1}^n K_{x_iu_i}\frac{\partial}{\partial \rho_{u_i}}  {S_\mu[\Psi](\rho_\mathbf{x},\bar{\rho}_\mathbf{x})} \bigg\lvert_{\rho_\mathbf{x}=0}
        \end{equation}
This expression allows us to recover the usual description of quantum states in the Literature.

\subsubsection{Quantizing operators}
At the same time, as the connection is also holomorphic, the quantization mapping $\mathcal{Q}$ is simple to compute and defines holomorphic linear operators on $\mathcal{H}_{Hol}$. In particular it is immediate to verify that
\begin{equation}
\label{eq:operator_field}
\mathcal{Q}(\phi^x) \Psi(\phi^\mathbf{x})=\phi^x \Psi(\phi^\mathbf{x}), \qquad \forall \Psi(\phi^\mathbf{x})\in \mathcal{H}_{Hol}
\end{equation}
\begin{equation}
    \label{eq:operator_momentum}
    \mathcal{Q}(\bar \phi^x) \Psi(\phi^\mathbf{x})=\Delta^{xy}\partial_{\phi^y }\Psi(\phi^\mathbf{x}),\  \forall \Psi(\phi^\mathbf{x})\in \mathcal{H}_{Hol}.
    \end{equation}

Notice that these operators correspond to the creation and annihilation operators, as it was to be expected  ${a^{\dagger x}}=\mathcal{Q}(\phi^x)$ while $a^x=\mathcal{Q}(\bar \phi^x)$.

\subsection{Schrödinger quantization}

Let us consider now the other main example, the case of Schrödinger quantization. Now, the quantum states must depend only on the field variables, and not on the field momenta.  This is achieved by a real polarization spanned by the momentum directions in $\mathcal{M}_C$. In any case, we must re-consider the construction above from a perspective based on real manifolds and not complex ones. 



\subsubsection{The measure and the connection}
Remember that we introduced the complex structure on the space $\mathcal{N}'\times \mathcal{N}'$, defined by the space of classical fields $ \varphi^\mathbf{x}$ and their corresponding field momenta $\pi^\mathbf{x}$. Once we introduce the complex structure, we can write that $\mathcal{N}'\times \mathcal{N}'\simeq \mathcal{N}'_\mathbb{C}$, although the identification depends on the particular complex structure introduced. Recall that, for the identification to be trivial we must use the set of coordinates $\tilde \varphi^\mathbf{x}$ and $\tilde \pi^\mathbf{x}$ of \eqref{eq:changeOfCoordinates}.  This complex space and its dual are introduced in the Gel'fand triple $\mathcal{N}_\mathbb{C}\subset \mathcal{H}_\Delta\subset \mathcal{N}_\mathbb{C}$ where $\mathcal{H}_\Delta$ is the Hilbert space defined by the one-particle states. Notice, though, that the spaces may still be treated as real spaces of twice the complex dimension, since the spaces are isomorphic. We will use, as above, the notation $\mathcal{M}_F$ to represent the real linear space defined by the Gel'fand triple and  $\mathcal{M}_C$ to represent the complex one. Both spaces are identical, as we argued above. The same applies to the corresponding tensors, in particular the symplectic forms $\omega_{\mathcal{M}_F}$ and $\omega_{\mathcal{M}_C}$, which represent the canonical symplectic form of $T^*\mathcal{N}'$ using the two  Gel'fand triples. And also to the line bundle $B$ required to define the prequantization: we construct identical bundles on the two identical base manifolds; and for each bundle, we can define a connection. These two connections need not to be identical (since we can consider several different connections for each case), but they must have a curvature proportional to the symplectic form. 

Identifying $\mathcal{M}_F$ and $\mathcal{M}_C$ also means that the measure on the space of fields can also be considered associated with the functional 
\begin{equation}
\label{eq:C_real}
C_S(\xi^{R}_\mathbf{x},\xi^I_\mathbf{x})=\int_{\mathcal{N}'\times \mathcal{N}'} D\mu(\tilde \varphi^\mathbf{x}, \tilde \pi^\mathbf{x}) e^{{i}\left(\xi_x^R\tilde \varphi^x+ \xi^I_x\tilde\pi^x\right)}.
\end{equation}
Considering $\rho_{\mathbf{x}}=\frac1{\sqrt{2}}(\xi^R_\mathbf{x}+i\xi^I_\mathbf{x})$ and $\phi^\mathbf{x}=\frac{1}{\sqrt2}(\tilde\varphi^{\mathbf{x}}+i \tilde\pi^{\mathbf{x}} )$, we must have that
\begin{equation}
\label{eq:C_equiv}
C_S(\xi^{R}_\mathbf{x},\xi^I_\mathbf{x})=C_H\left(\frac{\xi^{R}_\mathbf{x}+i\xi^{I}_\mathbf{x}}{2}, \frac{\xi^{R}_\mathbf{x}-i\xi^{I}_\mathbf{x}}{2}\right).
\end{equation}
 We conclude thus that the measure $D\mu$ is defined for the real-fields description, and coincides, up to an scale factor,  with the measure $D\mu_c$ constructed in the previous section in the complex description.

Again, we will consider a pre-quantum Hilbert space of sections of the bundle $B$, square-integrable with respect to the measure $D\mu$. We can also identify a vacuum section $\Psi_0^S$ compatible with the connection in the sense of Equation \eqref{eq:vacuum}. With this trivialization, the space of states becomes the set of square-integrable functions on $\mathcal{M}_F$ with respect to the measure $D\mu$. 

From the isomorphism discussed above between $\mathcal{M}_C$ and $\mathcal{M}_F$, we can also conclude the equivalence of the connections on both spaces and the corresponding covariant derivatives. Strictly speaking, we should use different notation for the vector fields with respect to which we differentiate, since they are isomorphic and not identical, but we will use the same, to make the manuscript easier to read.  We shall use then the notation of Equation \eqref{eq:covariantDerivative} to represent the covariant derivative in the real polarization case, keeping in mind the fact that the only relevant vector fields are those tangent to the submanifold spanned by the field directions, as we will see now.


\subsubsection{The real polarization}
In the Schrödinger case the polarization is spanned by the momentum directions, i.e. the real subspace $P_{S}=Im(H_\Delta)$. We can consider an adapted symplectic potential with respect to $P_S$ satisfying
\begin{equation}
\label{eq:symlectic_potential_S}
\theta_S(X)=0 \qquad \forall X\in P_S.
\end{equation}

This means that $\theta_S$ must be a linear combination of the set $\{ d\tilde \varphi^\mathbf{x} \}$ of "basical" forms. The coefficients are read from the identification of the Gel'fand triple and the canonical Liouville 1-form of $T^*\mathcal{N}'$. Writing it in terms of the variables we introduced in the previous section we obtain: 

\begin{equation}
\label{eq:symplecticAdaptedSch}
\theta_S=-{i}(\tilde\varphi^x+i\tilde\pi^x)K_{xy}d\tilde\varphi^y=\pi^x\delta_{xy}d\varphi^y-i\varphi^xK_{xy} d\varphi^y+\varphi^x(K A)_{xy}d\varphi^y.
\end{equation}
Notice that we added the term $-i\tilde\varphi^xK_{xy} d\tilde\varphi^y$ (which is just an exact one form equal to the differential of the norm of the $\bar\varphi^x$) to the usual conventions because in this way $2\Theta=\theta_S+\theta_M$ with \begin{equation}
\label{eq:symplecticmomentumnaive}
\theta_M= -(\tilde\varphi^x+i\tilde\pi^x)K_{xy}{d\tilde\pi^y}.
\end{equation} From this point of view $\theta_S$ is the restriction of $\Theta$ to $P_S$ directions with an scale factor of $2$ to keep up with the change of coordinates. In this way we can stablish the Schrödinger picture as a restriction of the holomorphic case to real directions. Nonetheless, this restriction must be taken with care. 

The first conclusion of this analysis on the polarization is that functions on $\mathcal{H}_{S}$ (i.e., the quantum states) must be functions depending on the submanifold $\mathcal{N}'$ with coordinates $\varphi^{\mathbf{x}}$ rather on the whole $\mathcal{M}_F$, as it was to be expected.  Indeed, as for any $X\in P_S$,  $\theta_S(X)=0$, 
\begin{equation}
\label{eq:condition_polari}
\nabla_X \Psi=0 \Rightarrow d\Psi(X)=0, \forall X\in P.
\end{equation}
As the vacuum state $\Psi_0^S$ must satisfy this condition, any other section in the set of square-integrable ones will be obtained as a product by a function which is also annihilated by the momentum directions. Nonetheless, the definition of a measure for this set of functions is not immediate, since the measure in the original scheme above was defined for functions depending on both the fields and their momenta. Nonetheless, we can repeat the construction for a different Gel'fand triple defined as
\begin{equation}
\label{eq:gelfand_sch}
\mathcal{N}\subset \mathcal{H}_\Delta^S \subset \mathcal{N}',
\end{equation}
where $\mathcal{H}_\Delta^S$ is the subspace of $\mathcal{H}_\Delta$ defined by the field states $ \xi^R_\mathbf{x}$. This triple defines a corresponding dual product $\langle \xi_\mathbf{x}, \varphi^\mathbf{x} \rangle= \xi_x\varphi^x$ which allows us to define a measure $D\mu_S$ by the functional
\begin{equation}
\label{eq:measure_sch}
C(\xi_{\mathbf{x}})=\int_{\mathcal{N}'} D\mu_S( \varphi^\mathbf{x}) e^{i \xi_x\varphi^x }.
\end{equation}
Notice that, by construction
\begin{equation}
\label{eq:measure_sch_C}
C(\xi_\mathbf{x})=C_S(\xi_\mathbf{x}, 0)= e^{-\frac14\xi_x\Delta^{xy}\xi_y}, 
\end{equation}
for $C_S$ the functional defined in Equation \eqref{eq:C_real}. Again, Minlos theorem ensures that there exists a unique measure satisfying this condition and therefore we can define the set of (polarized) quantum states to be the space of square-integrable functions (with respect to the vacuum state) $\mathcal{H}_S=L^2(\mathcal{N}', D\mu_S)$. 


\subsubsection{Quantizing operators}
We can readily quantize linear operators, such as the field operator $\varphi^x$ or the field momentum  $\pi^x$, whose action on the states $\Phi\in L^2(\mathcal{N}', D\mu_S)$ will be:
\begin{multline}
\label{eq:quantizationOfSchOperators}
\mathcal{Q}(\varphi^x) \Phi (\varphi^\mathbf{x})= \varphi^x \Phi(\varphi^\mathbf{x}),\\
 \mathcal{Q}(\pi^x\delta_{xy})\Phi (\varphi^\mathbf{x})=\left (-i\partial_{\varphi^y}+i{\varphi^zK_{zy}}-{\varphi^x(K A)_{xy}}\right ) \Phi(\varphi^\mathbf{x}).
\end{multline}

Analogously, we can also introduce creation and annihilation operators from the field and momentum operators as:

\begin{equation}
\label{eq:crationannihilation}
a^x= \frac{\Delta^{xy}\partial_{\varphi^y}-iA^x_{y}\varphi^y}{\sqrt{2}}, \qquad a^{\dagger x}=\sqrt{2}\varphi^x-\frac{\Delta^{xy}\partial_{\varphi^y}-iA^x_{y}\varphi^y}{\sqrt{2}}.
\end{equation}

These operators satisfy that  $\mathcal{Q}(\varphi^x)= \frac{a^x+a^{\dagger x}}{\sqrt2}$ and $\mathcal{Q}(\pi^x)= -iK^x_{y}\frac{a^y-a^{\dagger y}}{\sqrt2}$.
Notice that the complex term of \eqref{eq:symplecticAdaptedSch} is essential in order for the quantized momentum to be self adjoint  and it provides the right prescription for quantization   
 obtained in \cite{Corichi2004} from an algebraic rather than geometric quantization procedure. 
\subsection{Quantization in the Momentum Field space}
Another real polarization is $P_M=Re(\mathcal{H}_\Delta)$, in this case we could proceed as we have done so far with the Hilbert space associated to the measure $D\mu_c$ and the symplectic potential \eqref{eq:symplecticmomentumnaive} but it soon leads to cumbersome expressions for the quantized operators. This is because the change of coordinates \eqref{eq:changeOfCoordinates} is adapted to deal with the Schrödinger representation in the fields space, we will develop the Schrodinger picture in the momentum-field space. For now on we will denote it simply Momentum-field representation.  
\subsubsection{An adapted Hilbert space and Polarization }
 In order to quantize this theory in a way akin to the one adapted to the Schrödinger picture we should respect the momentum in the change of coordinates we modify  \eqref{eq:changeOfCoordinates}  to  
$$d{\breve{\varphi}^x}=(A^t)^x_yd\pi^y-D^x_yd{\varphi^x} \qquad{d\breve{\pi}^x}=d{\pi^x},$$
with $\phi^\mathbf{x}=\frac{1}{\sqrt2}(\breve\varphi^\mathbf{x}+i\breve\pi^\mathbf{x})$. This treatment, following the arguments below  \eqref{eq:changeOfCoordinates}, modifies \eqref{eq:C} and  the measure considered for this case is $D\nu_c$ defined by
\begin{equation}
    \label{eq:gaussianmeasuremomentum}
    \breve{C}(\rho_{\mathbf{x}}, \bar \rho_{\mathbf{x}})=\int_{\mathcal{M}_C}D\nu_c(\phi^\mathbf{x} )e^{i (\bar\rho_x\phi^x+\rho_x\bar \phi^x) }= e^{\bar\rho_xD^{xy}\rho_y }.
\end{equation}
This also leads to a different Holomorphic representation in $L^2_{Hol}(\mathcal{N}'_{\mathbb{C}},D\nu_c)$. 
Following our steps with the Schrödinger picture we define $D\nu_M$ as the measure choosing as characteristic functional for the momentum representation 

\begin{equation}
    \label{eq:measure_momentum_C}
    \breve{C}_M(\xi_\mathbf{x})= e^{\frac14\xi_xD^{xy}\xi_y}. 
    \end{equation}
Momentum polarization is therefore adapted to this Cameron Martin Hilbert space, $P_M=Re(\mathcal{H}_{-D})$ and we choose the adapted symplectic potential
\begin{equation}
\label{eq:symplectic}
\breve{\theta}_M=(\breve\varphi^x+i\breve\pi^x)D^{-1}_{xy}d\breve\pi^y=-\varphi^x\delta_{xy}d\pi^y+i\pi^xD^{-1}_{xy} d\pi^y+\pi^x(AD^{-1})_{xy}d\pi^y.
\end{equation} 
It follows from this prescription that 
\begin{multline}
    \label{eq:quantizationOfmomentumOperators}
    \mathcal{Q}(\pi^x) \Phi (\pi^\mathbf{x})= \pi^x \Phi(\pi^\mathbf{x}),\\
     \mathcal{Q}(\varphi^x\delta_{xy})\Phi (\pi^\mathbf{x})=\left (i\partial_{\pi^y}+i{\pi^zD^{-1}_{zy}}+{\pi^x(AD^{-1})_{xy}}\right ) \Phi(\pi^\mathbf{x}).
\end{multline}
In the momentum field space creation and annihilation operators are 
\begin{equation}
\label{eq:crationannihilationmomentum}
b^x= -\frac{D^{xy}\partial_{\pi^y}-i(A^t)^x_{y}\pi^y}{\sqrt{2}}, \qquad b^{\dagger x}=\sqrt{2}\pi^x+\frac{D^{xy}\partial_{\pi^y}-i(A^t)^x_{y}\pi^y}{\sqrt{2}}.
\end{equation}

These operate dually to creation and annihilation operators in the field space, they satisfy that  $\mathcal{Q}(\pi^x)= \frac{b^x+b^{\dagger x}}{\sqrt2}$ and $\mathcal{Q}(\varphi^x)= -i(D^{-1})^x_{y}\frac{b^y-b^{\dagger y}}{\sqrt2}$.

\subsection{Isomorphisms between Holomorphic, Schrödinger, and Momentum-field pictures.}
Once we have shown different representations of the same quantum field theory we must know to what extent they represent the same theory. As we will see, it turns out that we can relate every picture with unitary isomorphisms and therefore we can consider them to represent the same theory.   

\subsubsection{ Holomorphic and Schrodinger pictures: Segal-Bargmann modified transforms}
Slightly modifying the Segal-Bargmann transform defined in \autoref{def:Segaltranf} we can stablish an isomorphism $$\tilde{\mathcal{B}}_{Sch}:L^2(\mathcal{N}',D\mu_S)\to L^2_{Hol}(\mathcal{N}_{\mathbb{C}}', D\mu_c).$$  
In order to define this modified Segal-Bargmann transform we should deal with the extra $1/2$ factor appearing in the characteristic functional $C$ \eqref{eq:measure_sch_C}. To do so let us first define $\Psi_{Hol}=\tilde{\mathcal{B}}_{Sch}(\Psi_{Sch})$. Then, we choose
\begin{equation}
    \label{eq:segalBargmanntransform}
    \Psi_{Sch}(\varphi^{\mathbf{x}})=e^{{i}f(\varphi^\mathbf{x})}\int D\mu_S(\pi^{\mathbf{x}})\Psi_{Hol}\big(\sqrt{2}[\varphi^{\mathbf{x}}+i\pi^{\mathbf{x}}]\big),
\end{equation}
where $e^{if(\varphi^\mathbf{x})}$ is a phase factor. This is similar to the definition in \cite{Oeckl2012},  with the addition of the phase factor. We include it because the isomorphism does not respect the operator representation of the creation and annihilation operators and therefore modifies the representation of the momentum and field operators. For instance if we take $f=0$, the corresponding $\tilde{\mathcal{B}}_{0}$ transformation produces
$$\tilde{\mathcal{B}}_{0}^{-1}\phi^x \tilde{\mathcal{B}}_{0}=\sqrt{2}\varphi^x,\textrm{ and }\tilde{\mathcal{B}}_{0}^{-1}\partial_{\phi^x} \tilde{\mathcal{B}}_{0}= \frac{1}{\sqrt2}\partial_{\varphi^x}. $$
Thus, departing from the expressions of $a^x$ and $a^{\dagger x}$ shown under \eqref{eq:operator_momentum}  we can not recover the expressions of \eqref{eq:crationannihilation}. To solve this problem we choose $$ f(\varphi^\mathbf{x})=\frac{1}2\varphi^x(KA)_{xy}\varphi^y,$$ and obtain
\begin{gather}
\label{eq:beggman_seagal_sch_creat_anhil}
\tilde{\mathcal{B}}_{Sch}^{-1}\phi^x \tilde{\mathcal{B}}_{Sch}=\sqrt2\varphi^x-\frac{\Delta^{xy}\partial_{\varphi^y}-iA^x_{
    y}\varphi^y}{\sqrt 2},\nonumber\\ \tilde{\mathcal{B}}_{Sch}^{-1}\partial_{\phi^x} \tilde{\mathcal{B}}_{Sch}= \frac{\partial_{\varphi^x}-i(KA)_{x
y}\varphi^y}{\sqrt 2}.
\end{gather} 
With this choice we preserve the expected creation and annihilation expressions for each picture. This may be interpreted as a nontrivial phase in the relation of the Schrödinger and holomorphic vacua as is explained in \autoref{sec:Ambiguity}. Indeed, if $\Psi_0^{Hol}=1$ represents the vacuum of the Fock space in the holomorphic representation we obtain $\Psi^{Sch}_0=\exp(i\varphi^x(KA)_{xy}\varphi^y/2)$ as the representation of the vacuum in the Schrödinger representation.

Similarly for the momentum space we can stablish a unitary isomorphism with the antiholomorphic representation
$$\tilde{\overline{\mathcal{B}}}_{Mom}:L^2(\mathcal{N}',D\nu_M)\to L^2_{\overline{Hol}}(\mathcal{N}_{\mathbb{C}}', D\nu_c).$$
Let $\hat\Psi_{\overline{Hol}}=\tilde{\overline{\mathcal{B}}}_{Mom}(\hat\Psi_{Mom})$ then  
\begin{equation}
    \label{eq:segalBargmanntransformmomentum}
    \hat\Psi_{Mom}(\pi^{\mathbf{x}})=e^{\frac{i}2\pi^x(D^{-1}A)_{xy}\pi^y}\int D\nu_M(\varphi^{\mathbf{x}})\hat\Psi_{\overline{Hol}}\big(i\sqrt{2}[\varphi^{\mathbf{x}}-i\pi^{\mathbf{x}}]\big),
\end{equation}
thus it follows $\tilde{\overline{\mathcal{B}}}_{Mom}^{-1}\bar\phi^x \tilde{\overline{\mathcal{B}}}_{Mom}=ib^{\dagger x}$ and  $-\tilde{\overline{\mathcal{B}}}_{Mom}^{-1}D^{xy}\partial_{\bar\phi^y} \tilde{\overline{\mathcal{B}}}_{Mom}=-ib^x$.

\subsubsection{Schrodinger and Momentum-field representations: Fourier transforms. }
Holomorphic coordinates adapted to field or momenta representations are related via a holomorphic transformation 
\begin{equation}
\label{eq:changeofvar2}
(d\breve\varphi^x+id\breve\pi^x)=(K^x_y-iKA^x_y)(d\tilde\varphi^y+i d\tilde\pi^y).
\end{equation}
Using the relation, $D=(iA^t+\mathds{1})K(iA-\mathds{1})$, it is easy to see that
\begin{equation}
\label{eq:unitary}
\begin{array}{cccc}
U :&L^2_{Hol}(\mathcal{N}_{\mathbb{C}}', D\mu_c) & \longrightarrow & L^2_{\overline{Hol}}(\mathcal{N}_{\mathbb{C}}', D\nu_c) \\ 
 & \Psi_{Hol}(\phi^{\mathbf{x}}) &\longmapsto & \hat\Psi_{\overline{Hol}}(\bar\phi^{\mathbf{x}})=\Psi_{Hol}[i(D^{-1}-iAD^{-1})^x_y\bar\phi^y],  
\end{array}
\end{equation}
is an isometry relating holomorphic and anti-holomorphic functions,  with  inverse  
\begin{equation}
    \label{eq:unitaryinv}
    \begin{array}{cccc}
    U^{-1} :&L^2_{\overline{Hol}}(\mathcal{N}_{\mathbb{C}}', D\nu_c) & \longrightarrow & L^2_{Hol}(\mathcal{N}_{\mathbb{C}}', D\mu_c) \\ 
     & \hat\Psi_{\overline{Hol}}(\bar\phi^{\mathbf{x}}) &\longmapsto & \Psi_{Hol}(\phi^{\mathbf{x}})=\hat\Psi_{\overline{Hol}}[i(K+iKA)^x_y\phi^y].
    \end{array}
\end{equation}
Notice that the transformation acts non-trivially over creation and annihilation operators because $$U\phi^xU^{-1}=i(D^{-1}-iAD^{-1})^x_y\bar\phi^y$$
 and 
 $$U\partial_{\phi^x}U^{-1}=i(K+iKA)^y_x\partial_{\bar\phi^y}.$$


Using this mapping we are in disposition to define a Fourier transform 
\begin{equation}
\label{eq:Fourier_momentum}
\begin{array}{cccc}
\mathcal{F} :& L^2(\mathcal{N}',D\mu_S) & \longrightarrow & L^2(\mathcal{N}',D\nu_M) \\ 
 & \Psi(\varphi^{\mathbf{x}}) &\longmapsto & \hat\Psi(\pi^\mathbf{x}),  
\end{array}
\end{equation}
as
\begin{equation}
\label{eq:Fourier}
\mathcal{F} ={\overline{\mathcal{B}}}^{-1} U {\mathcal{B}}
\end{equation}
where ${\mathcal{B}}$ is the regular Segal-Bargmann transform and $\overline{\mathcal{B}}$ its antiholomorphic counterpart. We skip technicalities about factors $1/2$ in covariances that do not affect the discussion. Then, by construction, we get  
\begin{align}
\label{eq:Fourier_momentum_final}
\mathcal{F} \left(\sqrt{2}\varphi^x-\frac{\Delta^{xy}\partial_{\varphi^y}}{\sqrt2}\right)\mathcal{F}^{-1}&= i(\delta-iA )^x_y\left(\sqrt{2}D^{-1}\pi^y+\frac{\delta^{yz}\partial_{\pi^z}}{\sqrt{2}}\right) \nonumber \\
\mathcal{F} \Delta^{xy}\partial_{\varphi^y}\mathcal{F}^{-1}&= i(\delta+iA)^{xy}\partial_{\pi^y} 
\end{align}
from which follows  
\begin{align}
    \label{eq:momentum}
    \mathcal{F} \varphi^x\delta_{xy}\mathcal{F}^{-1}&=i\partial_{\pi^y}+i{\pi^zD^{-1}_{zy}}+{\pi^x(AD^{-1})_{xy}} \nonumber \\
    \mathcal{F}^{-1} \pi^{x}\delta_{xy}\mathcal{F}&= -i\partial_{\varphi^y}+i{\varphi^zK_{zy}}-{\varphi^x(K A)_{xy}}.
\end{align}
Therefore,  \eqref{eq:quantizationOfSchOperators} is transformed into  \eqref{eq:quantizationOfmomentumOperators} and viceversa.  Hence, the quantization mapping $\mathcal{Q}$ is respected by the Fourier transform. 

But this implies that creation and annihilation operators mix, in general,  in a nontrivial manner under $\mathcal{F}$:
\begin{align}
    \label{eq:Fourier_creation}
    \mathcal{F}a^x\mathcal{F}^{-1}&= i(\Delta-D^{-1})^x_y\frac{b^y}{2}+i(\Delta+D^{-1})^x_y\frac{b^{\dagger y}}{2} \nonumber \\
    \mathcal{F}a^{\dagger x}\mathcal{F}^{-1}&= -i(\Delta+D^{-1})^x_y\frac{b^y}{2}-i(\Delta-D^{-1})^x_y\frac{b^{\dagger y}}{2} 
    \end{align}
Notice that, from this expression, it is immediate that 
the mixing depends on the properties of the complex structure. Indeed, whenever $A^{\mathbf{x}}_\mathbf{y}=0$ we get $D^{-1}=-\Delta$. Hence, in those cases there will be no mixing between creation and annihilation operators under the Fourier transform.

\section{Quantization of higher order observables}
\label{sec:orderings}

In the geometric quantization framework described above, we have considered only classical linear observables because for more general functions, requirements on quantization  (axioms Q1-Q4 above)become more difficult to fulfill. Quantization mappings can not satisfy them in general, even for polynomial functions with degrees higher than 2. One possible way to circunvent this problem is to assume a dependence of the quantization mapping in $\hbar$ and ask axiom Q2) of the Dirac quantization conditions to hold only at leading order in $\hbar$. In this way, every quantization scheme must coincide for linear functions while it may differ for higher order polynomials. Thus, at it is well known, the different quantization mappings select a behavior with respect to the ordering problem. For instance, Weyl quantization will impose a symmetrized form on the polynomials in the field operators and their momenta,  while Wick or anti-Wick orders select the extrem form with respect to creation or annihilation operators. The non-trivial commutation relations of the field operator and their momenta, or, analogously, of the creation and annihilation operators, is the origin of the ordering problem, as we will see now  in some detail. We will work in the holomorphic representation of QFT because we know how to relate it with every other picture, as we explained above. 

\subsection{Weyl quantization}
As we just mentioned, 
when considering higher order observables me must deal with ordering problems of the $a^x,a^{\dagger y}$ operators, this is because of the nontrivial commutation relation $[a^x,a^{\dagger y}]=\Delta^{xy}$.  Weyl quantization mapping deals with this problem assigning to each monomial $\phi^n\bar\phi^m$ an averaged assignment of every possible order. We can then define the Weyl quantization to any polynomial classical functional inductively. For linear operators the quantization procedure must coincide with geometric quantization  $\mathcal{Q}_{Weyl}[\phi^x]=a^{\dagger x},   \mathcal{Q}_{Weyl}[\bar\phi^x]=a^{x}$. For higher order monomials, we proceed  symbolically as
\begin{multline}
\label{eq:weylaverage}
     \mathcal{Q}_{Weyl}[(\phi^n\bar\phi^m)^{u\vec{x},\vec{y}v}]=\\
   a^{\dagger u}\mathcal{Q}_{Weyl}[(\phi^{n-1}\bar\phi^m)^{\vec{x},\vec{y}v}]+\mathcal{Q}_{Weyl}[(\phi^{n-1}\bar\phi^m)^{\vec{x},\vec{y}v}]\frac{\Delta^{uz}}2\frac{\overleftarrow{\delta}}{\delta a^z}\\=\mathcal{Q}_{Weyl}[(\phi^{n}\bar\phi^{m-1})^{u\vec{x},\vec{y}}]a^{v}+\frac{\Delta^{zv}}2\frac{{\delta}}{\delta a^{\dagger z}}\mathcal{Q}_{Weyl}[(\phi^{n}\bar\phi^{m-1})^{u\vec{x},\vec{y}}]
\end{multline}
From this symbolic expression it follows that the action of Skorokhod integrals with covariance $\Delta^{\mathbf{xy}}/2$ is cast into
\begin{align}
\label{eq:WheylSkorokhod}
    \mathcal{Q}_{Weyl}[\partial^{*\phi^x}F]=\mathcal{Q}_{Weyl}[(\phi^x-{\frac{\Delta^{xy}}{2}\partial_{\bar\phi^y}})F]=a^{\dagger x}  \mathcal{Q}_{Weyl}[F]\nonumber \\
    \mathcal{Q}_{Weyl}[\partial^{*\bar\phi^x}F]=\mathcal{Q}_{Weyl}[(\bar\phi^x-{\frac{\Delta^{xy}}{2}\partial_{\phi^y}})F]=\mathcal{Q}_{Weyl}[F]a^{x}  
\end{align}


Thus this  quantization procedure is well suited to quantize classical functions belonging to a Hilbert space $\mathcal{H}_{cl}$ defined by
$$\mathcal{H}_{cl}=L^2(\mathcal{N}_{\mathbb{C}}, Dw)\textrm{ with }\int Dw(\phi^{\mathbf{x}})e^{i (\bar\rho_x\phi^x+\rho_x\bar \phi^x) }= e^{-\frac{\bar\rho_x\Delta^{xy}\rho_y}{2} }.$$
Indeed, for functions in this Hilbert space we can consider a chaos decomposition 

\begin{equation}
\label{eq:f_chaos}
F[\phi, \bar \phi] =\sum F^{(n,\bar{m})}_{\vec{x}\vec{y}}\mathcal{W}_{\frac{\Delta}{2}}(\phi^n\bar\phi^m)^{\vec{x}\vec{y}}.   
\end{equation}
Now, quantizing the functional just constructs the operator with the same chaos coefficients. For one such monomial we obtain
\begin{equation}
\label{eq:weyl}
\mathcal{Q}_{Weyl}\big[F^{(n,\bar{m})}_{\vec{x}\vec{y}}\mathcal{W}_{\frac{\Delta}{2}}(\phi^n\bar\phi^m)^{\vec{x}\vec{y}}\big]= F^{(n,\bar m)}_{\vec{x}\vec{y}}(a^{\dagger n})^{\vec{x}}(a^m)^{\vec{y}}.
 \end{equation}
It is straightforward to obtain the action on the total classical function by linearity.

Another property of Weyl quantization is the transformation complex conjugation into involution 
\begin{equation}
\label{eq:involution}
\mathcal{Q}_{Weyl}\big(\bar F\big)= \mathcal{Q}_{Weyl}\big(F\big)^{\dagger}.
\end{equation}
 Therefore, real classical functions  of $\mathcal{H}_{cl}$ are quantized into hermitian operators.

It  is important to notice that the resulting operator of \eqref{eq:weyl} is unbounded and therefore it fails  to be defined on the whole space of pure states $L^2_{Hol}(\mathcal{N}_{\mathbb{C}},D\mu_c)$. Moreover this procedure can be extended to more general classes of functions, such as Hida distributions, as long as the chaos decomposition in terms of $Dw$ exists. We will not consider now this case but we refer to  \cite{alonso2023a} for further details on these issues.

\subsection{Other orderings: Wick Quantization} 
The most common ordering in QFT is not Weyl quantization, but Wick quantization, which  assigns to regular monomials the so called \textit{normal order} prescription. From what we know from quantization of finite dimensional systems, we can expect that the corresponding quantization mapping will be related to the Weyl quantization  mapping by a suitable morphism of the resulting algebras. 

In order to find an explicit expression for this operation  we introduce  the Wick operator $\mathcal{W}_{\frac{\Delta}{2}}$ for general functionals. Its action on a monomial must re-order the terms to write the resulting expression in normal ordered. From the properties of the Skorokhod integrals \eqref{eq:holomorphicSkorokhod}, acting on the constant function 1, it is straightforward to obtain  the expression of this operator as 

\begin{equation}
    \mathcal{W}_{\frac{\Delta}{2}}=\exp\Big(-\frac{\Delta^{xy}\partial_{\bar\phi^x}\partial_{\phi^y}}{2}\Big). 
\label{eq:wick_iso} 
\end{equation}
Notice that we use the same symbol as for the case of the Wiener Ito chaos decompositon \eqref{eq:ComplexWienerItodec}, since it is the action of these operator on the suitable pair of fields what produces the suitable polynomials.  As an operator, it is invertible and
\begin{equation}
 \label{eq:wick-inverse}
 \mathcal{W}_{\frac{\Delta}{2}}^{-1}=\mathcal{W}_{-\frac{\Delta}2}.
\end{equation}

With this transformation, Wick quantization can be defined as
$$\mathcal{Q}_{Wick}= \mathcal{Q}_{Weyl}\circ\mathcal{W}_{\frac{\Delta}2},$$
where $\mathcal{Q}_{Weyl}$ is the mapping defined in the previous section.

In this way, it is immediate to verify that a quadratic Hamiltonian  as the one describing Klein Gordon, is quantized in Wick quantization as 
$$\mathcal{Q}_{Wick}(\phi^xC_{xy}\bar\phi^y)= a^{\dagger x}C_{xy}a^{y}.$$ 


\section{Conclusions}
\label{sec:conclusions}
In this paper we have been able to summarize the differential geometric aspects of the space of quantum states of a scalar field in terms of Hida-Malliavin calculus. At the level of the set of classical fields, we have learned to identify the role of the complex structure on the classical phase space of fields
\begin{itemize}
    \item to build the isomorphism of vectors and co-vectors on the nuclear space of test functions, required to build a classical model on the set of distributions on the spatial domain $\Sigma$ 
    \item to define the Hilbert space of quantum 1-particle states, and hence to use it as a reference for the Fock space of multiparticle states
\end{itemize}
We also discussed that the description of the theory in canonical or holomorphic coordinates changes the  Gel'fand triple structure that realizes the duality in the space of fields. This means that changes of coordinates affect hugely the analytical description of the model space of the manifold and those operations must be considered with care.  


Besides, we have characterized geometric quantization in this more complicated framework with different gauges for the polarization.  We have understood  the origin of the vacuum section as a reference and the relation of the vacuum phase and the symplectic potential defined on the classical phase space.  This systematic analysis led to the study of the Holomorphic, Schrödinger and momentum-field quantizations on an equal grounding. The tools of Malliavin calculus have been used to describe geometric quantization in such a way that we are now able to identify the meaning of  creation and annihilation operators in each space. This analysis also provides a systematic way of relating each picture with integral transforms, in this context we've been able to simplify, and expand the analytical expressions available in the literature for the Fourier and Segal-Bargmann transforms. This systematic approach has been useful to describe the role of the vacuum phase of the Schrödinger picture for each transform. We have also presented the momentum-field picture of QFT that, to our knowledge, has not been treated systematically in the literature.      

 We have also discussed the limitations of geometric quantization and we presented the Weyl quantization mapping for polynomial classical functionals using the Wiener-Ito chaos decomposition to define the representation of quantum operators on the (bosonic) Fock space in a straightforward way; and use it to define the Wick quantization mapping. 

 In the second part of this series we will present a different set of tools to study the Weyl and Wick quantization mappings in depth. We will also carry on with our geometrical description of the quantum phase space of pure states with the tools oh Hida calculus to describe a Poisson structure and finally introduce the time evolution and explore some examples.   



\textbf{Acknowledgments}
The authors would like to thank  Carlos Escudero Liébana for pointing out to very useful bibliography on the connection of Malliavin calculus with QFT.
The authors acknowledge partial financial support of Grant PID2021-
123251NB-I00 funded by MCIN/AEI/10.13039/501100011033 and by the
European Union, and of Grant E48-20R funded by Government of Aragón.
C.B-M and D.M-C acknowledge financial support by Gobierno de Aragón
through the grants defined in ORDEN IIU/1408/2018 and ORDEN
CUS/581/2020 respectively.




%\addcontentsline{toc}{chapter}{Bibliography}

\begin{thebibliography}{14}
    % BibTex style file: bmc-mathphys.bst (version 2.1), 2014-07-24
    \ifx \bisbn   \undefined \def \bisbn  #1{ISBN #1}\fi
    \ifx \binits  \undefined \def \binits#1{#1}\fi
    \ifx \bauthor  \undefined \def \bauthor#1{#1}\fi
    \ifx \batitle  \undefined \def \batitle#1{#1}\fi
    \ifx \bjtitle  \undefined \def \bjtitle#1{#1}\fi
    \ifx \bvolume  \undefined \def \bvolume#1{\textbf{#1}}\fi
    \ifx \byear  \undefined \def \byear#1{#1}\fi
    \ifx \bissue  \undefined \def \bissue#1{#1}\fi
    \ifx \bfpage  \undefined \def \bfpage#1{#1}\fi
    \ifx \blpage  \undefined \def \blpage #1{#1}\fi
    \ifx \burl  \undefined \def \burl#1{\textsf{#1}}\fi
    \ifx \doiurl  \undefined \def \doiurl#1{\url{https://doi.org/#1}}\fi
    \ifx \betal  \undefined \def \betal{\textit{et al.}}\fi
    \ifx \binstitute  \undefined \def \binstitute#1{#1}\fi
    \ifx \binstitutionaled  \undefined \def \binstitutionaled#1{#1}\fi
    \ifx \bctitle  \undefined \def \bctitle#1{#1}\fi
    \ifx \beditor  \undefined \def \beditor#1{#1}\fi
    \ifx \bpublisher  \undefined \def \bpublisher#1{#1}\fi
    \ifx \bbtitle  \undefined \def \bbtitle#1{#1}\fi
    \ifx \bedition  \undefined \def \bedition#1{#1}\fi
    \ifx \bseriesno  \undefined \def \bseriesno#1{#1}\fi
    \ifx \blocation  \undefined \def \blocation#1{#1}\fi
    \ifx \bsertitle  \undefined \def \bsertitle#1{#1}\fi
    \ifx \bsnm \undefined \def \bsnm#1{#1}\fi
    \ifx \bsuffix \undefined \def \bsuffix#1{#1}\fi
    \ifx \bparticle \undefined \def \bparticle#1{#1}\fi
    \ifx \barticle \undefined \def \barticle#1{#1}\fi
    \ifx \bconfdate \undefined \def \bconfdate #1{#1}\fi
    \ifx \botherref \undefined \def \botherref #1{#1}\fi
    \ifx \url \undefined \def \url#1{\textsf{#1}}\fi
    \ifx \bchapter \undefined \def \bchapter#1{#1}\fi
    \ifx \bbook \undefined \def \bbook#1{#1}\fi
    \ifx \bcomment \undefined \def \bcomment#1{#1}\fi
    \ifx \oauthor \undefined \def \oauthor#1{#1}\fi
    \ifx \citeauthoryear \undefined \def \citeauthoryear#1{#1}\fi
    \ifx \endbibitem  \undefined \def \endbibitem {}\fi
    \ifx \bconflocation  \undefined \def \bconflocation#1{#1}\fi
    \ifx \arxivurl  \undefined \def \arxivurl#1{\textsf{#1}}\fi
    \csname PreBibitemsHook\endcsname
    


     %%% 1
     \bibitem{Haag1964}
     \begin{barticle}
     \bauthor{\bsnm{Haag}, \binits{R.}},
     \bauthor{\bsnm{Kastler}, \binits{D.}}:
     \batitle{An {{Algebraic Approach}} to {{Quantum Field Theory}}}.
     \bjtitle{Journal of Mathematical Physics}
     \bvolume{5}(\bissue{7}),
     \bfpage{848}--\blpage{861}
     (\byear{1964}).
     \doiurl{10.1063/1.1704187}
     \end{barticle}
     \endbibitem
     
     %%% 2
     \bibitem{Streater1989}
     \begin{bbook}
     \bauthor{\bsnm{Streater}, \binits{R.F.}},
     \bauthor{\bsnm{Wightman}, \binits{A.S.}}:
     \bbtitle{{{PCT}}, {{Spin}} and {{Statistics}}, and {{All That}}}.
     \bpublisher{{Princeton University Press}},
     (\byear{1989}).
     \burl{https://www.jstor.org/stable/j.ctt1cx3vcq}
     Accessed 2023-03-18
     \end{bbook}
     \endbibitem
     
     %%% 3
     \bibitem{Doplicher1969}
     \begin{barticle}
     \bauthor{\bsnm{Doplicher}, \binits{S.}},
     \bauthor{\bsnm{Haag}, \binits{R.}},
     \bauthor{\bsnm{Roberts}, \binits{J.E.}}:
     \batitle{Fields, observables and gauge transformations {{I}}}.
     \bjtitle{Communications in Mathematical Physics}
     \bvolume{13}(\bissue{1}),
     \bfpage{1}--\blpage{23}
     (\byear{1969}).
     \doiurl{10.1007/BF01645267}
     \end{barticle}
     \endbibitem
     
     %%% 4
     \bibitem{Doplicher1969a}
     \begin{barticle}
     \bauthor{\bsnm{Doplicher}, \binits{S.}},
     \bauthor{\bsnm{Haag}, \binits{R.}},
     \bauthor{\bsnm{Roberts}, \binits{J.E.}}:
     \batitle{Fields, observables and gauge transformations. {{II}}}.
     \bjtitle{Communications in Mathematical Physics}
     \bvolume{15}(\bissue{3}),
     \bfpage{173}--\blpage{200}
     (\byear{1969}).
     \doiurl{10.1007/BF01645674}
     \end{barticle}
     \endbibitem
     
     %%% 5
     \bibitem{Doplicher1971}
     \begin{barticle}
     \bauthor{\bsnm{Doplicher}, \binits{S.}},
     \bauthor{\bsnm{Haag}, \binits{R.}},
     \bauthor{\bsnm{Roberts}, \binits{J.E.}}:
     \batitle{Local observables and particle statistics {{I}}}.
     \bjtitle{Communications in Mathematical Physics}
     \bvolume{23}(\bissue{3}),
     \bfpage{199}--\blpage{230}
     (\byear{1971}).
     \doiurl{10.1007/BF01877742}
     \end{barticle}
     \endbibitem
     
     %%% 6
     \bibitem{Doplicher1974}
     \begin{barticle}
     \bauthor{\bsnm{Doplicher}, \binits{S.}},
     \bauthor{\bsnm{Haag}, \binits{R.}},
     \bauthor{\bsnm{Roberts}, \binits{J.E.}}:
     \batitle{Local observables and particle statistics {{II}}}.
     \bjtitle{Communications in Mathematical Physics}
     \bvolume{35}(\bissue{1}),
     \bfpage{49}--\blpage{85}
     (\byear{1974}).
     \doiurl{10.1007/BF01646454}
     \end{barticle}
     \endbibitem
     
     %%% 7
     \bibitem{Buchholz1982}
     \begin{barticle}
     \bauthor{\bsnm{Buchholz}, \binits{D.}},
     \bauthor{\bsnm{Fredenhagen}, \binits{K.}}:
     \batitle{Locality and the structure of particle states}.
     \bjtitle{Communications in Mathematical Physics}
     \bvolume{84}(\bissue{1}),
     \bfpage{1}--\blpage{54}
     (\byear{1982}).
     \doiurl{10.1007/BF01208370}
     \end{barticle}
     \endbibitem
     
     %%% 8
     \bibitem{Buchholz2020}
     \begin{barticle}
     \bauthor{\bsnm{Buchholz}, \binits{D.}},
     \bauthor{\bsnm{Fredenhagen}, \binits{K.}}:
     \batitle{A {{C}}*-algebraic {{Approach}} to {{Interacting Quantum Field
       Theories}}}.
     \bjtitle{Communications in Mathematical Physics}
     \bvolume{377}(\bissue{2}),
     \bfpage{947}--\blpage{969}
     (\byear{2020}).
     \doiurl{10.1007/s00220-020-03700-9}
     \end{barticle}
     \endbibitem
     
     %%% 9
     \bibitem{connesNoncommutativeGeometry1994}
     \begin{bbook}
     \bauthor{\bsnm{Connes}, \binits{A.}}:
     \bbtitle{Noncommutative {{Geometry}}}.
     \bpublisher{{Academic Press}}
     (\byear{1994})
     \end{bbook}
     \endbibitem
     
     %%% 10
     \bibitem{connesRenormalizationQuantumField2000}
     \begin{barticle}
     \bauthor{\bsnm{Connes}, \binits{A.}},
     \bauthor{\bsnm{Kreimer}, \binits{D.}}:
     \batitle{Renormalization in {{Quantum Field Theory}} and the {{Riemann-Hilbert
       Problem I}}: {{The Hopf Algebra Structure}} of {{Graphs}} and the {{Main
       Theorem}}}.
     \bjtitle{Communications in Mathematical Physics}
     \bvolume{210}(\bissue{1}),
     \bfpage{249}--\blpage{273}
     (\byear{2000}).
     \doiurl{10.1007/s002200050779}
     \end{barticle}
     \endbibitem
     

    %%% 1
    \bibitem{Kibble1979}
    \begin{barticle}
    \bauthor{\bsnm{Kibble}, \binits{T.W.B.}}:
    \batitle{Geometrization of quantum mechanics}.
    \bjtitle{Communications in Mathematical Physics}
    \bvolume{65}(\bissue{2}),
    \bfpage{189}--\blpage{201}
    (\byear{1979}).
    \doiurl{10.1007/BF01225149}
    \end{barticle}
    \endbibitem
    
    %%% 2
    \bibitem{Hida1980}
    \begin{bbook}
    \bauthor{\bsnm{Hida}, \binits{T.}}:
    \bbtitle{Brownian {{Motion}}}.
    \bpublisher{{Springer US}},
    \blocation{{New York, NY}}
    (\byear{1980}).
    \doiurl{10.1007/978-1-4612-6030-1}
    \end{bbook}
    \endbibitem
    
    %%% 3
    \bibitem{Hida1993}
    \begin{bbook}
    \bauthor{\bsnm{Hida}, \binits{T.}},
    \bauthor{\bsnm{Kuo}, \binits{H.-H.}},
    \bauthor{\bsnm{Potthoff}, \binits{J.}},
    \bauthor{\bsnm{Streit}, \binits{L.}}:
    \bbtitle{White {{Noise}}}
    vol. \bseriesno{72}.
    \bpublisher{{Springer Netherlands}},
    \blocation{{Dordrecht}}
    (\byear{1993}).
    \doiurl{10.1007/978-94-017-3680-0}
    \end{bbook}
    \endbibitem
    
    %%% 4
    \bibitem{Wester2003}
    \begin{botherref}
    \oauthor{\bsnm{Westerkamp}, \binits{W.}}:
    Recent {{Results}} in {{Infinite Dimensional Analysis}} and {{Applications}} to
      {{Feynman Integrals}}
    (2003)
    {\href{https://arxiv.org/abs/math-ph/0302066}{{arXiv:math-ph/0302066}}}
    \end{botherref}
    \endbibitem
    
    %%% 5
    \bibitem{Corichi2004}
    \begin{barticle}
    \bauthor{\bsnm{Corichi}, \binits{A.}},
    \bauthor{\bsnm{Cortez}, \binits{J.}},
    \bauthor{\bsnm{Quevedo}, \binits{H.}}:
    \batitle{Schr\"odinger and {{Fock}} representation for a field theory on curved
      spacetime}.
    \bjtitle{Annals of Physics}
    \bvolume{313}(\bissue{2}),
    \bfpage{446}--\blpage{478}
    (\byear{2004}).
    \doiurl{10.1016/j.aop.2004.05.004}
    \end{barticle}
    \endbibitem
    
    %%% 6
    \bibitem{Oeckl2012}
    \begin{barticle}
    \bauthor{\bsnm{Oeckl}, \binits{R.}}:
    \batitle{The {{Schr\"odinger}} representation and its relation to the
      holomorphic representation in linear and affine field theory}.
    \bjtitle{Journal of Mathematical Physics}
    \bvolume{53}(\bissue{7}),
    \bfpage{072301}
    (\byear{2012})
    {\href{https://arxiv.org/abs/1109.5215}{{arXiv:1109.5215}}}.
    \doiurl{10.1063/1.4731770}
    \end{barticle}
    \endbibitem
    
    %%% 7
    \bibitem{Oeckl2012a}
    \begin{barticle}
    \bauthor{\bsnm{Oeckl}, \binits{R.}}:
    \batitle{Holomorphic quantization of linear field theory in the general
      boundary formulation}.
    \bjtitle{Symmetry, Integrability and Geometry: Methods and Applications
      (SIGMA)}
    \bvolume{8},
    \bfpage{31}
    (\byear{2012})
    {\href{https://arxiv.org/abs/1009.5615}{{arXiv:1009.5615}}}.
    \doiurl{10.3842/SIGMA.2012.050}
    \end{barticle}
    \endbibitem
    
    %%% 8
    \bibitem{Gellfand}
    \begin{botherref}
    \oauthor{\bsnm{Moiseevi{\v c}~Gel'fand}, \binits{I.}},
    \oauthor{\bsnm{Vilenkin}, \binits{N.Y.}}:
    Generalized functions: {{Applications}} of {{Harmonic Analysis}}, {{Volume}} 4
    (1967)
    \end{botherref}
    \endbibitem
    
    %%% 9
    \bibitem{Kriegl1997}
    \begin{bbook}
    \bauthor{\bsnm{Kriegl}, \binits{A.}},
    \bauthor{\bsnm{Michor}, \binits{P.}}:
    \bbtitle{The {{Convenient Setting}} of {{Global Analysis}}}
    vol. \bseriesno{53}.
    \bpublisher{{American Mathematical Society}},
    \blocation{{Providence, Rhode Island}}
    (\byear{1997}).
    \doiurl{10.1090/surv/053}
    \end{bbook}
    \endbibitem
    
    %%% 10
    \bibitem{Nunno2009}
    \begin{bbook}
    \bauthor{\bsnm{Nunno}, \binits{G.D.}},
    \bauthor{\bsnm{{\O}ksendal}, \binits{B.}},
    \bauthor{\bsnm{Proske}, \binits{F.}}:
    \bbtitle{Malliavin {{Calculus}} for {{L\'evy Processes}} with {{Applications}}
      to {{Finance}}}
    vol. \bseriesno{44}.
    \bpublisher{{Springer Berlin Heidelberg}},
    \blocation{{Berlin, Heidelberg}}
    (\byear{2009}).
    \doiurl{10.1007/978-3-540-78572-9}
    \end{bbook}
    \endbibitem
    
    %%% 11
    \bibitem{alonso2023a}
    \begin{botherref}
    \oauthor{\bsnm{Alonso}, \binits{J.L.}},
    \oauthor{\bsnm{{Bouthelier-Madre}}, \binits{C.}},
    \oauthor{\bsnm{{Clemente-Gallardo}}, \binits{J.}},
    \oauthor{\bsnm{{Mart{\'i}nez-Crespo}}, \binits{D.}}:
    Geometric Flavours of {{Quantum Field}} Theory on a {{Cauchy}} Hypersurface.
      {{Part II}}: {{Canonical}} and {{Geometrical QFT}}
    (2023)
    \end{botherref}
    \endbibitem
    
    %%% 12
    \bibitem{Brunetti2015}
    \begin{bbook}
    \bauthor{\bsnm{Brunetti}, \binits{R.}},
    \bauthor{\bsnm{Dappiaggi}, \binits{C.}},
    \bauthor{\bsnm{Fredenhagen}, \binits{K.}},
    \bauthor{\bsnm{Yngvason}, \binits{J.}}:
    \bbtitle{Advances in {{Algebraic Quantum Field Theory}}}.
    \bpublisher{{Springer International Publishing}},
    \blocation{{Cham}}
    (\byear{2015}).
    \doiurl{10.1007/978-3-319-21353-8}
    \end{bbook}
    \endbibitem
    
    %%% 13
    \bibitem{Woodhouse1997}
    \begin{bbook}
    \bauthor{\bsnm{Woodhouse}, \binits{N.M.J.}}:
    \bbtitle{Geometric {{Quantization}}},
    \bedition{Second edition} edn.
    \bsertitle{Oxford {{Mathematical Monographs}}}.
    \bpublisher{{Oxford University Press}},
    \blocation{{Oxford}}
    (\byear{1997})
    \end{bbook}
    \endbibitem
    
    %%% 14
    \bibitem{hall2013}
    \begin{bbook}
    \bauthor{\bsnm{Hall}, \binits{B.C.}}:
    \bbtitle{Quantum Theory for Mathematicians}.
    \bsertitle{Graduate Texts in Mathematics},
    vol. \bseriesno{267}.
    \bpublisher{{Springer}},
    \blocation{{New York}}
    (\byear{2013})
    \end{bbook}
    \endbibitem
    
    \end{thebibliography}


\end{document}
